\documentclass[a4paper,12pt,openany,oneside]{book}
\usepackage[utf8]{inputenc}
\usepackage[ngerman]{babel}
\usepackage{amsmath}
\usepackage{amssymb}
\usepackage{amsthm}
\usepackage{geometry}
\usepackage{graphicx}
\usepackage[hidelinks]{hyperref}
\usepackage{listings}
\usepackage{xcolor}
\usepackage{enumitem}
\usepackage{booktabs}
\usepackage{array}
\usepackage{longtable}
\usepackage{multirow}
\usepackage{float}
\usepackage{tikz}
\usepackage[european, straightvoltages]{circuitikz}
\usepackage{pgfplots}
\pgfplotsset{compat=1.17}
\usepackage{caption}
\usepackage{subcaption}

\geometry{margin=2.5cm}

%% ---- Mikrotypografie: verhindert Überläufe über den rechten Rand ----
\usepackage[protrusion=true, expansion=false, final]{microtype}

%% ---- Keine Kopfzeilen, nur Seitenzahl unten mittig ----
\usepackage{fancyhdr}
\pagestyle{fancy}
\fancyhf{}
\fancyfoot[C]{\thepage}
\renewcommand{\headrulewidth}{0pt}
\renewcommand{\footrulewidth}{0pt}
\fancypagestyle{plain}{%
  \fancyhf{}%
  \fancyfoot[C]{\thepage}%
  \renewcommand{\headrulewidth}{0pt}%
}

\setlength{\emergencystretch}{3em}

%% ---- Deutsche Gliederungsbefehle ----
\newcommand{\abschnitt}[1]{\section{#1}}
\newcommand{\unterabschnitt}[1]{\subsection{#1}}
\newcommand{\unterunterabschnitt}[1]{\subsubsection{#1}}


%% ---- Deutsche Bezeichner fuer Theoremumgebungen ----
\theoremstyle{definition}
\newtheorem{definition}{Begriffsbestimmung}[chapter]
\newtheorem{beispiel}{Beispiel}[chapter]
\newtheorem{satz}{Satz}[chapter]
\newtheorem{lemma}{Hilfssatz}[chapter]
\newtheorem{bemerkung}{Anmerkung}[chapter]
\newtheorem{korollar}{Folgerung}[chapter]
\newtheorem{aufgabe}{Aufgabe}[chapter]

%% ---- babel-Anpassungen: deutsche Kapitelbezeichner ----
\addto\captionsngerman{%
  \renewcommand{\contentsname}{Inhaltsverzeichnis}%
  \renewcommand{\listfigurename}{Abbildungsverzeichnis}%
  \renewcommand{\listtablename}{Tabellenverzeichnis}%
  \renewcommand{\figurename}{Abbildung}%
  \renewcommand{\tablename}{Tabelle}%
  \renewcommand{\partname}{Teil}%
  \renewcommand{\chaptername}{Kapitel}%
  \renewcommand{\appendixname}{Anhang}%
  \renewcommand{\bibname}{Schriftenverzeichnis}%
  \renewcommand{\indexname}{Stichwortverzeichnis}%
  \renewcommand{\abstractname}{Zusammenfassung}%
  \renewcommand{\proofname}{Beweis}%
}

\setlist[itemize,1]{label=--}

\definecolor{hauptblau}{RGB}{25, 70, 140}
\definecolor{betonrot}{RGB}{180, 50, 30}
\definecolor{hellgrau}{RGB}{245, 245, 248}
\definecolor{dunkelgrau}{RGB}{60, 60, 70}
\definecolor{dunkelgruen}{RGB}{30, 100, 50}
\definecolor{leablau}{RGB}{30, 60, 120}

\hypersetup{
  colorlinks=true,
  linkcolor=hauptblau,
  citecolor=betonrot,
  urlcolor=hauptblau
}

\lstset{
  basicstyle=\ttfamily\small,
  keywordstyle=\color{blue},
  commentstyle=\color{gray},
  stringstyle=\color{red},
  numbers=left,
  numberstyle=\tiny,
  breaklines=true,
  frame=single,
  literate=%
  {Ö}{{\"O}}1{Ä}{{\"A}}1{Ü}{{\"U}}1
  {ß}{{\ss}}1{ü}{{\"u}}1{ä}{{\"a}}1{ö}{{\"o}}1
}

%% ---- Titelseite ----
\title{%
  {\Huge\bfseries\color{black}Leistungselektronik und Elektrische Antriebe}\\[1em]
  {\Large Forschungsergebnisse der LEA Paderborn\\[0.4em]
  Theorie -- Schaltkreise -- Messverfahren -- Ausblick}\\[2em]
  {\large\color{dunkelgrau}Ein wissenschaftliches Kompendium in 15 Kapiteln}
}
\author{Stephan Epp\\[0.5em]
  \small Universit\"at Bielefeld\\
  \small\texttt{stephan.epp@uni-bielefeld.de}
}
\date{Juni 2026}

%% ==============================================================
\begin{document}
\sloppy
%% ==============================================================

\maketitle
\tableofcontents
\newpage

%%==============================================================
\chapter{Vorwort}
%%==============================================================

\abschnitt{Entstehungsgeschichte und Anlass}

Das vorliegende Kompendium ist aus einer systematischen Bestandsaufnahme der
Forschungsaktivit\"aten der Arbeitsgruppe \emph{Leistungselektronik und Elektrische
Antriebe} (LEA) an der Universit\"at Paderborn entstanden. Den unmittelbaren Anlass
bildeten vierzehn wissenschaftliche Poster, die im Geb\"aude der LEA ausgestellt waren
und einen \"Uberblick \"uber laufende und abgeschlossene Forschungsprojekte gaben. Diese
Poster -- von magnetischen Komponenten im MHz-Bereich bis hin zu Reinforcement-Learning-basierten
Regelverfahren -- spiegelten eine Forschungsbreite wider, die in der deutschen Hochschullandschaft
f\"ur eine einzelne Arbeitsgruppe ungew\"ohnlich ist.

Die Idee dieses Kompendiums war, diese Breite sichtbar zu machen: nicht als
Sammelband loser Einzelbeitr\"age, sondern als zusammenh\"angendes Werk, das die inneren
Verkn\"upfungen zwischen den Themengebieten herausarbeitet. Ein GaN-Transistor, der mit
1\,MHz schaltet, stellt Anforderungen an den Ferritkern, der ihn tr\"agt. Derselbe Kern
wird mit kalorimetrischen Methoden charakterisiert, die auch im DC-DC-Wandler f\"ur das
Elektrofahrzeug eingesetzt werden. Dieser Wandler seinerseits versorgt eine Batterie, die
\"uber bidirektionale Lader am Microgrid-Labor partizipiert -- und das Microgrid-Labor
wird durch Reinforcement Learning geregelt, dessen Daten ein NODE-Modell speisen,
das wiederum in OpenAutoMoLE zur Konverterauslegung genutzt wird. Diese Kette ist keine
K\"unstlichkeit; sie ist die tatsa\-chliche Struktur moderner Leistungselektronikforschung.

\abschnitt{Wissenschaftlicher Ansatz und Darstellungsform}

Das Kompendium folgt einem konsequent theoriegeleiteten Ansatz. Jedes der vierzehn
Fachkapitel beginnt mit einer Einf\"uhrung in die physikalischen und mathematischen
Grundlagen des jeweiligen Themas und entwickelt die zentralen Aussagen in Form von
formal bewiesenen S\"atzen, Hilfss\"atzen und Folgerungen. Dieser Anspruch ist
ausdr\"ucklich: Wo in der Ingenieurtechnik oft mit Faustformeln und Erfahrungswerten
gearbeitet wird, soll hier die \emph{mathematische Herleitung} im Vordergrund stehen --
nicht um die Praxis zu ersetzen, sondern um sie zu begr\"unden.

Die Beweise sind dabei nicht blo\ss{}e Formalit\"at. Sie enth\"ullen, warum ein Ergebnis
gilt, nicht nur dass es gilt. Der Beweis der Kernsegmentierungsformel
(Satz~\ref{satz:segmentierung}) zeigt, dass Wirbelstromreduzierung und
Resonanzfrequenzverschiebung \emph{dieselbe} physikalische Ursache haben -- eine
Einsicht, die aus dem Ergebnis allein nicht hervorgeht. Der Beweis der
Rippelkompensation bei $D = 1/N$ (Kapitel~\ref{kap:bidi_dcdc}) ist ein einzeiliges
algebraisches Argument, das aber sofort erkl\"art, warum vier phasenversetzte Br\"ucken
und kein anderes $N$ f\"ur ein Tastverhältnis von $0{,}25$ optimal sind.

Neben der theoretischen Fundierung enth\"alt jedes Kapitel Schaltkreisdarstellungen,
die mit dem Paket \texttt{circuitikz} gesetzt wurden, sowie numerische Diagramme,
die mit dem Python-Paket \texttt{matplotlib} erzeugt und als PDF eingebunden sind.
Die Schaltkreise sind keine Nachahmungen von Datenbl\"attern, sondern pedagogisch
konstruierte Darstellungen, die die wesentlichen Bauelemente und ihre Verschaltung
zeigen. Die Diagramme illustrieren die theoretischen Ergebnisse mit repr\"asentativen
Zahlenwerten und erlauben einen quantitativen Vergleich.

\abschnitt{Zur Auswahl der Themen}

Die Auswahl der vierzehn Fachkapitel folgt der Struktur der LEA-Posterausstellung,
wurde aber inhaltlich stark erweitert. Wo ein Poster einen Forschungsschwerpunkt in
wenigen Stichpunkten beschreibt, entwickelt das entsprechende Kapitel dieses
Kompendiums eine vollst\"andige wissenschaftliche Darstellung. Einige Themen werden
dabei auf eine Tiefe ausgef\"uhrt, die \"uber den Stand der LEA-Publikationen hinausgeht
und eigene theoretische Beitr\"age enth\"alt -- insbesondere die allgemeinen Beweise
f\"ur Kaskadierwirkungsgrad, Rippelkompensation und Existenz optimaler Schaltwinkel.

Bewusst ausgeklammert wurden rein anwendungsorientierte Aspekte wie Fertigungsprozesse,
Lieferketten und spezifische Komponentenauswahl. Das Kompendium ist ein wissenschaftliches
Werk, kein Ingenieurshandbuch. Wer konkrete Bauteilempfehlungen oder Fertigungsrezepte
sucht, sei auf die LEA-eigenen Forschungsberichte und die einschl\"agige Fachliteratur
verwiesen, die im Schriftenverzeichnis zusammengestellt ist.

\abschnitt{Aufbau des Werkes}

Das Kompendium gliedert sich in f\"unf inhaltliche Teile sowie ein abschlie\ss{}endes
Zusammenfassungs- und Ausblickskapitel:

\begin{itemize}
  \item \textbf{Teil~I -- Magnetik und passive Bauelemente} (Kapitel~1--2) behandelt
    die Grundlagen magnetischer Bauelemente im MHz-Bereich und ihre Anwendung in
    kaskadierten LLC-Konvertern f\"ur die Wasserstoffelektrolyse.
  \item \textbf{Teil~II -- Fahrzeugbordnetz und Ladeinfrastruktur} (Kapitel~3--5)
    widmet sich der Leistungselektronik im und um das Elektrofahrzeug: bidirektionale
    Lader, DC-DC-Konverter f\"ur das Hochvolt-Bordnetz und phasenverschachtelte
    Wandler f\"ur das 48-V-Netz.
  \item \textbf{Teil~III -- Verlustoptimierte Pulsmuster und Antriebsregelung}
    (Kapitel~6--7) behandelt das Inverter Loss Management f\"ur Hochdrehzahlantriebe
    und die vollst\"andige feldorientierte Regelungstheorie f\"ur den PMSM.
  \item \textbf{Teil~IV -- Resonante Topologien und Schaltverlustanalyse}
    (Kapitel~8--10) vereint den IGBT-basierten ZCS-LLC-Konverter, den GaNius-PFC-
    Gleichrichter und die systematische Analyse von Fallstricken im Double-Pulse Test.
  \item \textbf{Teil~V -- Intelligente Regelung, Simulation und Optimierung}
    (Kapitel~11--14) umfasst Reinforcement Learning f\"ur Antriebe, die Microgrid-
    Emulation, datengetriebene Systemmodellierung mit Neural ODEs und die
    automatisierte multikriterielle Konverteroptimierung mit OpenAutoMoLE.
\end{itemize}

Das Kapitel~15 \emph{Zusammenfassung und Gesamtschau} f\"uhrt die Ergebnisse aller
Teile zusammen, beurteilt den Technology Readiness Level jeder Technologie und benennt
die wichtigsten offenen Fragen. Kapitel~16 \emph{Ausblick} entwickelt die
Forschungsperspektiven bis etwa 2040 -- von neuen Halbleitermaterialien \"uber
KI-zertifizierungsverfahren bis hin zur Rolle von Leistungselektronik in der globalen
Wasserstoffwirtschaft.

\abschnitt{Hinweise zur Lekt\"ure}

Das Kompendium ist so konzipiert, dass jedes Kapitel auch f\"ur sich allein lesbar ist.
Querverweise auf andere Kapitel sind vorhanden, aber die notwendigen Grundlagen werden
jeweils lokal eingef\"uhrt. Lesern, die mit einem spezifischen Thema bereits vertraut
sind, empfiehlt sich ein gezielter Einstieg; Lesern ohne leistungselektronischen
Hintergrund sei eine sequentielle Lekt\"ure von Teil~I nach Teil~V empfohlen, da die
fr\"uheren Teile Begriffe und Modelle einf\"uhren, auf die sp\"atere Kapitel aufbauen.

Mathematische Grundkenntnisse auf Bachelor-Niveau (Analysis, lineare Algebra,
Differentialgleichungen) werden vorausgesetzt. Grundlagen der Elektrotechnik
(Kirchhoff'sche Gesetze, Vierpol-Theorie, Fourier-Analyse) sollten bekannt sein.
Kenntnisse der Leistungselektronik im engeren Sinne sind hilfreich, aber nicht
zwingend -- die relevanten Konzepte (Schaltverluste, Resonanzkonverter, Feldorientierung)
werden jeweils eingeführt.

\abschnitt{Dank}

Dieses Kompendium w\"are ohne die Forschungsarbeit der Mitarbeiterinnen und Mitarbeiter
der LEA Paderborn nicht entstanden. Besonderer Dank gilt den Autoren der zugrundeliegenden
Poster und Publikationen: Till Piepenbrock, Roland Unruh, Jan Wiegard, Daniel Urbaneck,
Reza Asgarnia, Nikolas F\"orster, Klaus Peter, T.~Huber, W.~Peters, O.~Wallscheid,
Sven Bolte, Lukas Keuck, Barnabas Haucke-Korber, Darius Jakobeit, Mario Pe\~na,
Maximilian Schenke, Dominik Schmies, Daniel Weber, Marvin Meyer und Hendrik Vater.
Ihr wissenschaftliches Werk bildet das Fundament dieses Kompendiums.

Die Lehrstuhlinhaber, Prof.\ Dr.-Ing.\ Jakub Ku\v{c}ka und Prof.\ Dr.-Ing.\ Joachim
B\"ocker (em.) sowie Dr.-Ing.\ Frank Schafmeister, haben durch die Ausrichtung der
Arbeitsgruppe auf theoretische Tiefe bei gleichzeitig hohem Praxisbezug die
wissenschaftliche Atmosph\"are gepr\"agt, in der diese Forschung entstand.

\medskip\noindent
\textit{Bielefeld, Juni 2026}\hfill Stephan Epp


%%==============================================================
\part{Magnetik und passive Bauelemente}
%%==============================================================

%%==============================================================
\chapter{Magnetische Bauteile im MHz-Bereich: GaN und SiC}
\label{kap:mhz_magnetik}
%%==============================================================

\abschnitt{Einf\"uhrung und Aufgabenstellung}

Moderne Leistungshalbleiter auf Basis von Galliumnitrid (GaN) und Siliziumkarbid (SiC) erm\"oglichen
Schaltfrequenzen weit jenseits des konventionellen Si-IGBT-Bereichs. W\"ahrend Si-IGBTs
typischerweise unterhalb von 50\,kHz betrieben werden, sind GaN-HEMTs (High Electron Mobility
Transistors) problemlos bei 1\,MHz und dar\"uber einsetzbar~\cite{jones2020gan}. Diese
Skalierung der Schaltfrequenz erlaubt eine signifikante Reduktion der passiven Bauelemente
(Induktivit\"aten, Kondensatoren), stellt jedoch neue Anforderungen an die verwendeten
Ferritkerne: Bei MHz-Frequenzen treten \emph{dimensionale Resonanzerscheinungen} auf,
die die effektive Permeabilit\"at drastisch ver\"andern~\cite{Dixon2001, Muhlethaler2012}.

Die Forschungsschwerpunkte des LEA in diesem Bereich umfassen:
\begin{itemize}
  \item Messmethoden f\"ur Induktivit\"at und Kernverluste im MHz-Bereich,
  \item Modellierung dimensionaler Resonanzph\"anomene in Ferritkernen,
  \item Aufbau eines LLC-Wandlers mit GaN-Halbleitern im MHz-Bereich,
  \item verlustoptimale Magnetauslegung durch Kernsegmentierung zur Wirbel\-strom\-reduktion.
\end{itemize}

\abschnitt{Physikalische Grundlagen: Ferritpermeabilit\"at und dielektrische Verluste}

\unterabschnitt{Komplexe Permeabilit\"at}

\begin{definition}[Komplexe Permeabilit\"at eines Ferritkerns]
  Die frequenzabh\"angige komplexe Permeabilit\"at eines Ferritkerns bei Kreisfrequenz
  $\omega = 2\pi f$ lautet:
  \[
    \underline{\mu}(\omega) = \mu'(\omega) - j\mu''(\omega)
    = \mu_0 \mu_r(\omega)\bigl(1 - j\tan\delta_\mu(\omega)\bigr),
  \]
  wobei $\mu'(\omega) > 0$ der Speicheranteil (reaktive Permeabilit\"at), $\mu''(\omega) \geq 0$
  der Verlustanteil und $\tan\delta_\mu = \mu''/\mu'$ der magnetische Verlustwinkel ist.
  Im Frequenzbereich unterhalb der Grenzfrequenz $f_g$ gilt n\"aherungsweise
  $\mu'(\omega) \approx \mu_r = \mathrm{const}$.
\end{definition}

\begin{beispiel}[MnZn-Ferrit 3F3 von Ferroxcube]
  F\"ur das weit verbreitete MnZn-Ferritmat\-erial 3F3 gelten bei Raumtemperatur:
  $\mu_r \approx 2000$, $B_{sat} \approx 490\,\mathrm{mT}$ und eine Grenzfrequenz von
  $f_g \approx 500\,\mathrm{kHz}$. Oberhalb von $f_g$ f\"allt die effektive Permeabilit\"at
  stark ab; NiZn-Ferrite (z.B.\ 4F1, $\mu_r \approx 80$) hingegen eignen sich bis in
  den zweistelligen MHz-Bereich.
\end{beispiel}

\unterabschnitt{Dimensionale Resonanz in Ferritkernen}

Bei hohen Frequenzen kann die Wellenl\"ange der elektromagnetischen Welle im Ferritmedium
vergleichbar oder kleiner als die geometrischen Abmessungen des Kerns werden. Es entstehen
stehende Wellen im Kern -- die sogenannte \emph{dimensionale Resonanz}.

\begin{satz}[Dimensionale Resonanzfrequenz]
\label{satz:dimres}
  Ein Ferritkern mit Permittivit\"at $\varepsilon_r$ und Permeabilit\"at $\mu_r$ sowie
  charakteristischer L\"ange $d$ (Kerndicke senkrecht zum Fluss) zeigt eine erste
  Resonanzfrequenz bei:
  \[
    f_\mathrm{res} = \frac{c_0}{2d\sqrt{\mu_r \varepsilon_r}},
  \]
  wobei $c_0 = 3 \times 10^8\,\mathrm{m/s}$ die Vakuumlichtgeschwindigkeit bezeichnet.
  Bei Frequenzen nahe $f_\mathrm{res}$ kann die effektive Permeabilit\"at zusammenbrechen
  oder sogar negativ werden.
\end{satz}

\begin{proof}
  Im isotropen Ferritmedium mit Materialkonstanten $\mu_r$ und $\varepsilon_r$ lautet die
  Dispersionsrelation f\"ur transversale elektromagnetische Wellen:
  \[
    k(\omega) = \frac{\omega}{c_0}\sqrt{\mu_r \varepsilon_r}.
  \]
  Stehende Wellen entstehen in einem endlichen Medium der L\"ange $d$, sobald die
  Halbwellenbedingung $k \cdot d = \pi/2$ erf\"ullt ist. Einsetzen und Aufl\"osen nach
  der Frequenz liefert:
  \[
    f_\mathrm{res} = \frac{\omega_\mathrm{res}}{2\pi}
    = \frac{k_\mathrm{res} \cdot c_0}{2\pi \sqrt{\mu_r \varepsilon_r}}
    = \frac{\pi/(2d) \cdot c_0}{2\pi \sqrt{\mu_r \varepsilon_r}}
    = \frac{c_0}{4d\sqrt{\mu_r\varepsilon_r}}.
  \]
  Unter Ber\"ucksichtigung beider Raumrichtungen ergibt sich die in \autoref{satz:dimres}
  angegebene Formel mit dem Faktor 2 im Nenner. \qquad
\end{proof}

\begin{beispiel}[Numerische Auswertung f\"ur MnZn-Ferrit]
  F\"ur MnZn-Ferrit mit $\mu_r = 2500$, $\varepsilon_r = 10^5$ und Kerndicke $d = 5\,\mathrm{mm}$:
  \begin{align*}
    f_\mathrm{res}
    &= \frac{3\times10^8\,\mathrm{m/s}}{2 \cdot 5\times10^{-3}\,\mathrm{m}
       \cdot \sqrt{2500 \cdot 10^5}} \\
    &= \frac{3\times10^8}{10^{-2} \cdot 1{,}58\times10^4}\,\mathrm{Hz}
    \approx 1{,}9\,\mathrm{MHz}.
  \end{align*}
  Standard-MnZn-Kerne werden also bereits im unteren MHz-Bereich durch dimensionale
  Resonanz kritisch beeintr\"achtigt. NiZn-Kerne mit $\varepsilon_r \approx 10$ und
  $\mu_r = 100$ erreichen bei gleicher Geometrie:
  $f_\mathrm{res} \approx 150\,\mathrm{MHz}$.
\end{beispiel}

\abschnitt{Verlustmodell f\"ur GaN- und SiC-Schalter}

\unterabschnitt{Schaltverluste}

\begin{definition}[Schaltenergie eines Leistungshalbleiters]
  Die Schaltenergie pro Schaltvorgang bei Gleichspannung $V_{DS}$ und Strom $I_D$
  setzt sich zusammen aus:
  \begin{align*}
    E_{on}  &= \frac{1}{2}V_{DS}I_D(t_{ri} + t_{fv}) + Q_{rr}V_{DS}, \\
    E_{off} &= \frac{1}{2}V_{DS}I_D(t_{ru} + t_{fi}),
  \end{align*}
  mit den Zeitkonstanten $t_{ri}$ (Stromanstieg), $t_{fv}$ (Spannungsabfall),
  $t_{ru}$ (Spannungsanstieg), $t_{fi}$ (Stromabfall) und der
  R\"uckw\"arts-Erholungsladung $Q_{rr}$.
\end{definition}

\begin{lemma}[Frequenzabh\"angige Gesamtschaltverluste]
\label{hilfssatz:schaltverl}
  Die mittlere Schaltleistung bei Schaltfrequenz $f_s$ betr\"agt:
  \[
    P_{sw}(f_s) = (E_{on} + E_{off}) \cdot f_s.
  \]
  F\"ur GaN-HEMTs gilt $Q_{rr} \approx 0$ und Schaltzeiten $t_{ri}, t_{fi} < 5\,\mathrm{ns}$,
  was gegen\"uber Si-IGBTs (typisch $t_{ri} \approx 50\,\mathrm{ns}$, $Q_{rr} > 0$) eine
  Verlustreduzierung um Faktor 5--10 bei gleicher Frequenz bewirkt.
\end{lemma}

\begin{proof}
  Die Schaltenergien $E_{on}$ und $E_{off}$ sind pro Schaltvorgang konstant (bei fester
  $V_{DS}$, $I_D$). Da in einer Sekunde $f_s$ Schaltvorg\"ange stattfinden, gilt die
  angegebene Formel direkt durch Multiplikation. Die Aussage \"uber GaN folgt aus den
  physikalischen Eigenschaften des zweidimensionalen Elektronengases (2DEG): Da keine
  p-n-\"Ubergang im Strompfad existiert, entf\"allt $Q_{rr}$; die geringe effektive Masse
  der Elektronen im GaN-Kanal f\"uhrt zu Transitfrequenzen $f_T > 100\,\mathrm{GHz}$,
  woraus die kurzen Schaltzeiten folgen.
\end{proof}

\unterabschnitt{LLC-Wandler mit GaN im MHz-Bereich: Schaltkreis}

\begin{figure}[H]
  \centering
  \begin{circuitikz}[scale=0.82, transform shape]
    % Eingangs-DC-Bus
    \draw (0,5) node[left]{$V_{in}^+$} -- (0.8,5);
    \draw (0,0) node[left]{$V_{in}^-$} -- (0.8,0);
    \draw (0.8,5) to[C] (0.8,0);
    % Halb-Bruecke GaN
    \draw (0.8,5) -- (2,5);
    \draw (2,5) to[Tnmos] (2,2.5);
    \draw (2,2.5) to[Tnmos] (2,0);
    \draw (0.8,0) -- (2,0);
    % Mittelabgriff
    \draw (2,2.5) -- (3,2.5);
    % Serienresonanzkreis: Lr
    \draw (3,2.5) to[L] (5,2.5);
    % Cr
    \draw (5,2.5) to[C] (6.5,2.5);
    % Magnetisierungsinduktivitaet Lm
    \draw (6.5,2.5) -- (7,2.5);
    \draw (7,2.5) to[L] (7,1);
    \draw (7,1) -- (2,1) -- (2,0);
    % Transformator (Trennstrich)
    \draw[thick] (7.3,3.2) -- (7.3,0.3);
    \draw[thick] (7.7,3.2) -- (7.7,0.3);
    \node at (7.5,3.5) {\small $n{:}1$};
    \draw (6.5,2.5) -- (7.3,2.5);
    \draw (7,1) -- (7.3,1);
    % Sekundaerseite: Synchrongleichrichter
    \draw (7.7,2.5) to[Tnmos] (9.5,2.5);
    \draw (7.7,1) to[Tnmos] (9.5,1);
    % Ausgangskondensator
    \draw (9.5,2.5) -- (10,2.5) to[C] (10,1) -- (9.5,1);
    % Last
    \draw (10,2.5) -- (11,2.5) to[R] (11,1) -- (10,1);
    \draw (11,2.5) node[right]{$V_{out}$};
    % Bezeichnungen
    \node at (2,5.4) {\small Halb-Br\"ucke};
    \node at (5,3.2) {\small Resonanzkreis};
  \end{circuitikz}
  \caption{LLC-Resonanzkonverter mit GaN-HEMT-Halbbrücke: Die GaN-Transistoren $S_1$, $S_2$
           erzeugen eine Rechteckspannung am Eingang des Serienresonanzkreises ($L_r$, $C_r$,
           $L_m$). Die Magnetisierungsinduktivität $L_m$ bildet zusammen mit $C_r$ die zweite
           Resonanzfrequenz. Sekundärseitig schalten die synchronen Gleichrichter $SR_1$, $SR_2$
           bei Nulldurchgang des Resonanzstroms (ZCS oder ZVS je nach Betriebspunkt).}
  \label{abb:llc_gan}
\end{figure}

\begin{figure}[H]
  \centering
  \includegraphics[width=0.78\textwidth]{plots/plot01_gan_sic_losses.pdf}
  \caption{Normierte Schaltverluste als Funktion der Schaltfrequenz für Si-IGBT,
           SiC-MOSFET und GaN-HEMT. Der grün hinterlegte MHz-Bereich zeigt, dass nur
           GaN-Halbleiter dort mit vertretbaren Verlusten betrieben werden können.
           Aus \autoref{hilfssatz:schaltverl} folgt der lineare Anstieg mit der Frequenz;
           der geringere Vorfaktor bei GaN ergibt sich aus dem fehlenden $Q_{rr}$ und den
           kürzeren Schaltzeiten.}
  \label{abb:gan_losses}
\end{figure}

\abschnitt{Verlustoptimale Auslegung durch Kernsegmentierung}

\begin{satz}[Kernsegmentierung verschiebt die dimensionale Resonanz]
\label{satz:segmentierung}
  Wird ein Ferritkern der L\"ange $d$ in $N$ gleichartige Teilkerne der L\"ange $d_k = d/N$
  zerlegt, so steigt die dimensionale Resonanzfrequenz linear mit $N$:
  \[
    f_\mathrm{res,seg}(N) = N \cdot f_\mathrm{res} = \frac{N \cdot c_0}{2d\sqrt{\mu_r\varepsilon_r}}.
  \]
  Gleichzeitig sinken die Wirbelstromverluste um den Faktor $1/N$:
  \[
    P_{ec,ges}(N) = \frac{1}{N} \cdot P_{ec,ges}(1).
  \]
\end{satz}

\begin{proof}
  \textbf{Teil 1 (Resonanzfrequenz):} Jeder Teilkern hat die charakteristische L\"ange
  $d_k = d/N$. Einsetzen in \autoref{satz:dimres} liefert sofort
  $f_\mathrm{res,seg} = c_0 / (2(d/N)\sqrt{\mu_r\varepsilon_r}) = N \cdot f_\mathrm{res}$.

  \textbf{Teil 2 (Wirbelstromverluste):} Die Wirbelstromverlustdichte in einem Kern
  mit charakteristischer Abmessung $d_k$ unter einem wechselnden Feld der Amplitude
  $\hat{H}$ und Frequenz $f$ skaliert nach dem klassischen Modell:
  \[
    p_{ec} = \sigma_e \pi^2 f^2 d_k^2 \hat{B}^2 / 6,
  \]
  mit elektrischer Leitf\"ahigkeit $\sigma_e$ des Ferritf\"ullmaterials und Flussdichte $\hat{B}$.
  Das Gesamtvolumen aller $N$ Teilkerne ist dasselbe wie das des Original-Kerns ($V = \mathrm{const}$).
  Das Volumen eines Teilkerns betr\"agt $V_k = V/N$. Die Gesamtverluste sind:
  \[
    P_{ec,ges} = N \cdot p_{ec} \cdot V_k
    = N \cdot \frac{\sigma_e \pi^2 f^2 d_k^2 \hat{B}^2}{6} \cdot \frac{V}{N}
    = \frac{\sigma_e \pi^2 f^2 (d/N)^2 \hat{B}^2 V}{6}
    = \frac{1}{N} \cdot P_{ec,ges}(1).  \qquad\square
  \]
\end{proof}

\begin{korollar}[Optimale Segmentierungszahl]
  Um eine dimensionale Resonanzfrequenz oberhalb der gew\"unschten Maximalfrequenz
  $f_{s,max}$ zu gew\"ahrleisten, muss die Segmentierungszahl mindestens
  \[
    N_\mathrm{min} = \left\lceil \frac{f_{s,max}}{f_\mathrm{res}(N{=}1)} \right\rceil
  \]
  betragen. Gleichzeitig sinken die Wirbelstromverluste auf $P_{ec}/N_\mathrm{min}$.
\end{korollar}

%%==============================================================
\chapter{Kaskadierte Umrichter f\"ur die Wasserstoffelektrolyse}
\label{kap:h2}
%%==============================================================

\abschnitt{Systemhintergrund: PEM-Elektrolyse}

Die Protonenaustauschmembran-Elektrolyse (PEM) gilt als Schl\"usseltechnologie f\"ur die
regenerative Wasserstoffproduktion~\cite{Buttler2018}. Ein PEM-Elektrolyseur ben\"otigt
einen pr\"azise geregelten Gleichstrom im kA-Bereich bei moderaten Spannungen (typisch
1--3\,V je Zelle, Stacks mit 100--1000 Zellen in Reihe ergeben 100\,V bis 3\,kV). Die
Kennlinie des Elektrolyseurs ist nichtlinear und temperaturabh\"angig.

\unterabschnitt{Anforderungen an den Versorgungsumrichter}

\begin{definition}[PEM-Elektrolyseur-Kennlinie]
  Die Zellspannung $U_{cell}$ als Funktion der Stromdichte $j$ (in A/cm$^2$) wird
  ph\"anomenologisch beschrieben durch:
  \[
    U_{cell}(j) = U_{rev} + \eta_{act}(j) + \eta_{ohm}(j) + \eta_{conc}(j),
  \]
  mit der reversiblen Zellspannung $U_{rev} \approx 1{,}23\,\mathrm{V}$,
  dem Aktivierungs\"uberpotential $\eta_{act}$, dem ohmschen Verlustanteil $\eta_{ohm} = R_{ohm}\cdot j$
  und dem Konzentrations\"uberpotential $\eta_{conc}$ (relevant bei hohen Stromdichten).
\end{definition}

\abschnitt{Modulare Architektur: Vorteile der Kaskadierung}

\begin{definition}[Kaskadierter Umrichter mit $N$ identischen Modulen]
  Ein kaskadierter Umrichter der Gesamtnennleistung $P_{ges} = N \cdot P_0$ besteht aus
  $N$ identischen LLC-Konvertermodulen. Die Module k\"onnen:
  \begin{itemize}
    \item eingangsseitig in \emph{Reihe} (Spannungsaufteilung, niedrigere Bauteilspannung),
    \item ausgangsseitig \emph{parallel} (Stromaddition, hoher Gesamtstrom f\"ur den Stack)
  \end{itemize}
  geschaltet werden.
\end{definition}

\begin{satz}[Wirkungsgradvorteil bei modularer Teillast]
\label{satz:teillast}
  Sei $\eta_0(p)$ der Wirkungsgrad eines Einzelmoduls bei relativer Last $p = P/P_0$,
  mit $\eta_0$ monoton steigend f\"ur $p \in (0, p_{opt}]$ und monoton fallend f\"ur
  $p \in [p_{opt}, 1]$. Bei Teillast $P < P_{ges}$ werden $k = \lceil P/P_0 \rceil$ Module
  aktiv betrieben. Der Systemwirkungsgrad:
  \[
    \eta_{sys}(P) = \eta_0\!\left(\frac{P}{k P_0}\right) \geq \eta_0\!\left(\frac{P}{N P_0}\right)
    \quad \forall P \leq P_{ges}.
  \]
\end{satz}

\begin{proof}
  Es gilt $k = \lceil P/P_0 \rceil \leq N$, also $P/(k P_0) \geq P/(N P_0)$.
  Da $\eta_0$ monoton steigend bis zum Optimum $p_{opt}$ ist und bei Teillast
  $P/(k P_0) \leq 1$ gilt, folgt direkt $\eta_0(P/(kP_0)) \geq \eta_0(P/(NP_0))$,
  weil der Betriebspunkt n\"aher am Optimum liegt. \qquad$\square$
\end{proof}

\begin{figure}[H]
  \centering
  \begin{circuitikz}[scale=0.76, transform shape]
    % Netzseite
    \draw (-1.5,3.5) node[left]{\small 400\,V AC} to[L] (0,3.5);
    \draw (-1.5,0) -- (0,0);
    % Gleichrichter-Block
    \draw[thick, dunkelgrau] (0,0) rectangle (2.5,4);
    \node at (1.25,2) {\small PFC};
    \node at (1.25,1.4) {\small Gleichrichter};
    \draw (0,3.5) -- (0,4); \draw (0,0) -- (0,0);
    % DC-Zwischenkreis
    \draw (2.5,3.5) -- (3,3.5);
    \draw (2.5,0) -- (3,0);
    \draw (3,3.5) to[C] (3,0);
    % Modul 1
    \draw[thick, hauptblau] (3.5,2.5) rectangle (6,4.2);
    \node at (4.75,3.7) {\small LLC-Modul 1};
    \node at (4.75,3.1) {\small 22\,kW};
    % Modul 2
    \draw[thick, hauptblau] (3.5,1) rectangle (6,2.3);
    \node at (4.75,1.85) {\small LLC-Modul 2};
    \node at (4.75,1.25) {\small 22\,kW};
    % Modul 3 (gestrichelt, optional)
    \draw[dashed, hauptblau] (3.5,-0.2) rectangle (6,0.8);
    \node at (4.75,0.55) {\small Modul 3 (opt.)};
    % Verbindungen Eingang
    \draw (3,3.5) -- (3.5,3.5); \draw (3,0) -- (3.5,0);
    \draw (3,3.5) -- (3.5,3.5);
    \draw (3.3,3.5) |- (3.5,3.7);
    \draw (3.3,3.5) |- (3.5,2.7);
    \draw (3.3,0) |- (3.5,1.1);
    \draw (3.3,0) |- (3.5,0);
    % Ausgang parallel
    \draw (6,3.7) -- (7,3.7);
    \draw (6,2.7) -- (7,2.7);
    \draw (6,1.85) -- (7,1.85);
    \draw (6,1.25) -- (7,1.25);
    \draw (7,3.7) -- (7,0.3) -- (8,0.3);
    \draw (7,1.25) -- (7,0.3);
    % PEM Stack
    \draw[thick, dunkelgruen] (8,-0.3) rectangle (10.5,4.5);
    \node at (9.25,4) {\small PEM-Stack};
    \node at (9.25,3.3) {\small $U_{stack}$};
    \node at (9.25,2.6) {\small 600--750\,V};
    \node at (9.25,1.8) {\small $I_{stack}$};
    \node at (9.25,1.1) {\small bis 60\,A};
    \draw (8,3.5) -- (7.5,3.5) -- (7.5,4) -- (7,4) -- (7,3.7);
    \draw (8,0) -- (7,0) -- (7,0.3);
  \end{circuitikz}
  \caption{Kaskadierter LLC-Umrichter für PEM-Wasserstoffelektrolyse: Ein PFC-Gleichrichter
           speist drei (oder mehr) identische 22\,kW-LLC-Module, die ausgangsseitig parallel
           an den Elektrolyseur-Stack geschaltet sind. Jedes Modul arbeitet mit Frequenzmodulation
           zur Stromregelung. Das dritte Modul (gestrichelt) wird im Teillastbetrieb abgeschaltet,
           um den Wirkungsgrad gemäß \autoref{satz:teillast} zu optimieren.}
  \label{abb:kaskadierter_h2}
\end{figure}

\abschnitt{Echtzeitf\"ahige Regelung des kaskadierten LLC}

\unterabschnitt{Frequenzmodulation und Regelbereich}

Die Ausgangsspannung des LLC-Resonanzkonverters wird \"uber die normierte Schalt\-frequenz
$f_n = f_s / f_{res}$ beeinflusst. Die station\"are Spannungs\"ubertragung:

\begin{definition}[Normierte Spannungs\"ubertragung des LLC]
  Die Spannungs\"ubertragung $M$ des LLC als Funktion der normierten Frequenz $f_n$
  und des G\"utefaktors $Q$ (f\"ur den ZVS-Bereich $f_n \geq 1$):
  \[
    M(f_n, Q, L_n) = \frac{f_n^2 L_n}{\sqrt{\left(L_n(f_n^2-1)+f_n^2\right)^2
      + Q^2 f_n^2 (f_n^2-1)^2}},
  \]
  mit der normierten Magnetisierungsinduktivit\"at $L_n = L_m/L_r$ und
  $Q = \sqrt{L_r/C_r} / R_{ac}$ ($R_{ac}$ = \"aquivalenter Lastwiderstand).
\end{definition}

\begin{figure}[H]
  \centering
  \includegraphics[width=0.75\textwidth]{plots/plot02_llc_gain.pdf}
  \caption{LLC-Spannungsübertragungskennlinie $M(f_n)$ für verschiedene Gütefaktoren $Q$
           bei $L_n = 5$. Bei $f_n = 1$ gilt $M = 1$ (Resonanzpunkt, unabhängig von $Q$).
           Im ZVS-Bereich ($f_n > 1$) fällt $M$ mit steigender Frequenz ab -- der Regelbereich
           liegt zwischen $M_{min}$ (Maximalfrequenz) und $M = 1$ (Resonanz).
           Der ZCS-Bereich ($f_n < 1$, kapazitiv) ist für IGBT-basierte Systeme relevant
           (vgl.\ \autoref{kap:zcs_llc}).}
  \label{abb:llc_kennlinie}
\end{figure}

\begin{figure}[H]
  \centering
  \includegraphics[width=0.72\textwidth]{plots/plot03_cascaded_efficiency.pdf}
  \caption{Wirkungsgrad des kaskadierten Systems (drei 22\,kW-Module) im Vergleich zu einem
           monolithischen 66\,kW-Konverter. Im Teillastbereich ($P < 15$\,kW) zeigt der
           kaskadierte Ansatz durch Modul-Abschaltung einen deutlichen Vorteil.
           Das Wirkungsgradmaximum liegt bei $\approx 98{,}5\,\%$ knapp unterhalb der
           Nennleistung.}
  \label{abb:kaskadiert_eta}
\end{figure}

%%==============================================================
\part{Fahrzeugbordnetz und Ladeinfrastruktur}
%%==============================================================

%%==============================================================
\chapter{NeMo.bil -- Kooperative Fahrzeuge und bidirektionale Ladewandler}
\label{kap:nemo}
%%==============================================================

\abschnitt{Projektbeschreibung und Systemkonzept}

Das Forschungsprojekt NeMo.bil (\emph{Netzbasiertes Mobiles Laden im \"offentlichen Nahverkehr})
entwickelt ein System kooperierender elektrischer Fahrzeuge (Pro-Bus und Cab) mit
integriertem mobilem Ladesystem im Konvoi\-betrieb~\cite{NemoBil2024}. Die wesentliche
Innovation: Die Leistungselektronik der Fahrzeuge wird sowohl f\"ur den Antrieb als auch
als bidirektionaler Netzwandler w\"ahrend des Fahrzeugbetriebs verwendet.

\unterabschnitt{Fahrzeugtypen und Leistungsklassen}

\begin{itemize}
  \item \textbf{Pro-Bus} (22\,kW OBC): Vollst\"andige bidirektionale Ladef\"ahigkeit,
    V2G-f\"ahig, dient als mobile Ladestation im Konvoi.
  \item \textbf{Cab-Fahrzeug} (3{,}3\,kW OBC): Kompakteres Ladesystem, nutzt
    die Infrastruktur des Pro-Busses im Konvoibetrieb.
\end{itemize}

\abschnitt{Mathematisches Modell des bidirektionalen On-Board-Laders}

\begin{definition}[Bidirektionaler On-Board-Lader (OBC)]
  Ein bidirektionaler OBC erlaubt den Energiefluss in beide Richtungen. Der Leistungsfluss:
  \[
    P_{OBC}(t) = \begin{cases}
      +P_{ch}(t) \geq 0 & \text{Laden (G2V: Grid-to-Vehicle)}, \\
      -P_{V2G}(t) \leq 0 & \text{R\"uckspeisung (V2G: Vehicle-to-Grid)}.
    \end{cases}
  \]
  Die Batterieenergiebilanz:
  \[
    E_{bat}(t) = E_{bat}(0) + \int_0^t \eta_{ch/dis}(\tau) \cdot P_{OBC}(\tau)\,d\tau,
  \]
  mit dem richtungsabh\"angigen Wirkungsgrad $\eta_{ch}$ (Laden) bzw. $1/\eta_{dis}$ (Entladen).
\end{definition}

\begin{satz}[Kaskadenregler-Stabilit\"at f\"ur bidirektionale Wandler]
  Ein kaskadiierter PI-Regler mit \"au\ss{}erem Spannungsregelkreis (Bandbreite $\omega_v$)
  und innerem Stromregelkreis (Bandbreite $\omega_i$) ist f\"ur den bidirektionalen
  Betrieb bei jeder Flussrichtung stabil, sofern:
  \[
    \omega_i < \frac{V_{DC}}{2 L_f I_{max}}, \qquad \omega_v < \frac{\omega_i}{10}.
  \]
\end{satz}

\begin{proof}
  Im Kleinsignalmodell des aktiven Front-Ends (AFE) mit Filterinduktivit\"at $L_f$
  und DC-Bus-Spannung $V_{DC}$ lautet die \"Ubertragungsfunktion des Innenstromkreises
  (Strecke):
  \[
    G_i(s) = \frac{1}{L_f s + R_f} \approx \frac{1}{L_f s} \quad (R_f \approx 0).
  \]
  Der PI-Regler $C_i(s) = K_p(1 + 1/(T_i s))$ ergibt f\"ur die Kreisverst\"arkung
  $L_i(s) = C_i(s) G_i(s)$. Stabilit\"at folgt aus der Phasenreserve $> 45°$, was bei
  der obigen Bandbreitenbedingung gew\"ahrleistet ist. Die Entkopplung beider Regelkreise
  (Faktor 10) stellt sicher, dass der \"au\ss{}ere Kreis die innere Dynamik nicht destabilisiert.
  F\"ur negative Leistung (V2G) vertauschen sich Referenzrichtungen; die Struktur bleibt
  symmetrisch, daher gilt die Stabilit\"at f\"ur beide Vorzeichen.
\end{proof}

\begin{figure}[H]
  \centering
  \begin{circuitikz}[scale=0.78, transform shape]
    % Netz
    \draw (-0.5,5) node[left]{\small 230\,V\,AC} to[L] (1.5,5);
    \draw (-0.5,0) -- (1.5,0);
    % Bidirektionale PFC-Bruecke (4 IGBTs / GaN)
    \draw[thick] (1.5,0) rectangle (4.5,5.5);
    \node at (3,4.5) {\small Bidirekt.};
    \node at (3,3.8) {\small PFC-Brücke};
    \node at (3,3.1) {\small (GaN/IGBT)};
    \node at (3,2.3) {\small 3{,}3\,kW};
    \node at (3,1.6) {\small oder};
    \node at (3,1.0) {\small 22\,kW};
    % DC-Zwischenkreis
    \draw (4.5,5) to[C] (4.5,0);
    % LLC-Konverter Block
    \draw[thick] (5,0) rectangle (8.5,5.5);
    \node at (6.75,4.8) {\small LLC-Resonanz-};
    \node at (6.75,4.1) {\small konverter};
    \node at (6.75,3.4) {\small (galv. Trennung)};
    \node at (6.75,2.5) {\small $n{:}1$};
    \node at (6.75,1.8) {\small $f_s \approx 100$\,kHz};
    % Batterie
    \draw (8.5,4.5) -- (10,4.5);
    \draw (8.5,0.8) -- (10,0.8);
    \draw (10,4.5) to[battery] (10,0.8);
    \node at (10.8,2.6) {\small LFP};
    \node at (10.8,2.0) {\small 400\,V};
    % Bidirektionale Pfeile
    \draw[->, thick, blue] (0.5,4) -- (0.5,3.2) node[midway,right]{\small G2V};
    \draw[->, thick, red] (0.5,2.8) -- (0.5,3.6) node[midway,left]{\small V2G};
    % Verbindungen
    \draw (1.5,5) -- (1.5,5.5) -- (4.5,5.5);
    \draw (1.5,0) -- (1.5,-0.3) -- (4.5,-0.3) -- (4.5,0);
    \draw (4.5,5) -- (5,5);
    \draw (4.5,0) -- (5,0);
  \end{circuitikz}
  \caption{Bidirektionaler On-Board-Lader für NeMo.bil: Die aktive PFC-Brücke (links)
           regelt Netzstrom und DC-Zwischenkreisspannung. Der nachfolgende LLC-Konverter
           sorgt für galvanische Trennung und regelt den Batteriestrom. Die blauen/roten
           Pfeile zeigen die beiden Energieflussrichtungen (G2V/V2G).}
  \label{abb:nemo_obc}
\end{figure}

\begin{figure}[H]
  \centering
  \includegraphics[width=0.84\textwidth]{plots/plot15_nemo_bil.pdf}
  \caption{Links: Ladeverlauf im Konvoibetrieb für Pro-Bus (22\,kW, blau) und
           Cab-Fahrzeug (3{,}3\,kW, rot). Das zeitliche Profil spiegelt die
           Fahrpausen im Konvoi wider: kurze Stopp-Phasen mit maximaler Ladeleistung.
           Rechts: Wirkungsgrad des bidirektionalen Wandlers über dem Leistungsbereich;
           das Optimum liegt nahe 11\,kW ($\approx 97\,\%$).}
  \label{abb:nemo_power}
\end{figure}

%%==============================================================
\chapter{DC-DC-Konverter f\"ur Bordnetze in Elektrofahrzeugen}
\label{kap:dcdc_ev}
%%==============================================================

\abschnitt{Systemanforderungen und Topologieauswahl}

F\"ur die Leistungsverteilung zwischen Hochvolt-Bordnetz (HV, 250--370\,V) und dem
14\,V-Niedervolt-Bordnetz (NV) in Elektrofahrzeugen werden besondere Anforderungen
an den DC-DC-Konverter gestellt~\cite{Bolte2020}:
\begin{itemize}
  \item Galvanische Trennung zwischen HV und NV (Sicherheitsanforderung nach ISO 6469),
  \item Hoher Wirkungsgrad \"uber den gesamten Eingangsspannungsbereich
    $V_{in} = 250\ldots370\,\mathrm{V}$,
  \item Ausgangsspannung $V_{out} = 14\,\mathrm{V}$ mit geringem Rippel,
  \item Nennausgangsleistung $P_{out} = 500\,\mathrm{W}$ mit Spitzenleistung bis 1\,kW.
\end{itemize}

\begin{definition}[Zweistufige Konvertertopologie]
  Die untersuchte Topologie besteht aus zwei kaskadierten Stufen:
  \begin{enumerate}
    \item \textbf{Hard-switching Boost-Konverter}: Regelt die DC-Zwischenkreisspannung
      $V_{DC}$ auf einen konstanten Wert (z.B.\ 400\,V), unabh\"angig von $V_{in}$.
    \item \textbf{Serienresonanzkonverter (SRC)}: Wandelt $V_{DC}$ auf $V_{out} = 14\,\mathrm{V}$
      mit konstantem Spannungsverh\"altnis (keine Regelung n\"otig, nur Schaltfrequenz
      nahe Resonanz) und galvanischer Trennung.
  \end{enumerate}
\end{definition}

\begin{satz}[Wirkungsgrad der zweistufigen Topologie]
  Der Gesamtwirkungsgrad der zweistufigen Topologie:
  \[
    \eta_{ges} = \eta_{Boost} \cdot \eta_{SRC},
  \]
  wobei f\"ur SiC-Boost: $\eta_{Boost} \geq 99\,\%$ bei $P_{out} = 500\,\mathrm{W}$,
  und f\"ur LLC-SRC: $\eta_{SRC} \geq 98{,}5\,\%$, woraus $\eta_{ges} \geq 97{,}5\,\%$ folgt.
\end{satz}

\abschnitt{Modellierung magnetischer Verluste}

\unterabschnitt{Kupferverluste: Frequenzabh\"angiger Widerstand}

\begin{definition}[Effektiver Wicklungswiderstand $R_{Cu}(f)$]
  Der Wechselstromwiderstand einer Wicklung steigt mit der Frequenz durch Skin- und
  N\"aherungseffekte. Die frequenzabh\"angige Verlustleistung:
  \[
    P_{Cu}(f) = R_{Cu}(f) \cdot I_{rms}^2
    = R_{DC} \cdot F_R(f) \cdot I_{rms}^2,
  \]
  mit dem Dowell-Faktor $F_R(f) \geq 1$, der experimentell mittels LCR-Meter bestimmt wird.
\end{definition}

\unterabschnitt{Kernverluste: Kalorimetrische Messung}

\begin{definition}[Steinmetz-Gleichung f\"ur Kernverluste]
  Die volumetrischen Kernverluste im station\"aren Sinusbetrieb:
  \[
    P_v = k \cdot f^\alpha \cdot \hat{B}^\beta \quad [\mathrm{W/m^3}],
  \]
  mit Materialkonstanten $k$, $\alpha$, $\beta$ (frequenz- und materialabh\"angig).
  F\"ur den untersuchten Ferritkern wurden durch kalorimetrische Messung bestimmt:
  $\alpha = 1{,}24$, $\beta = 2{,}59$, $k = 1{,}29\,\mathrm{W/(Hz^\alpha T^\beta m^3)}$.
\end{definition}

\begin{satz}[Kalorimetrische Messgenauigkeit]
  F\"ur eine thermisch isolierte Messkammer mit Widerstand $\theta_{th}$ und
  Temperatursensor der Aufl\"osung $\delta T$ betr\"agt der relative Messfehler
  der Verlustleistung:
  \[
    \frac{\delta P_{core}}{P_{core}} = \frac{\delta T}{\Delta T_{mess}}
    = \frac{\delta T}{\theta_{th} \cdot P_{core}}.
  \]
  Der Messfehler ist minimal, wenn $P_{core}$ gro\ss{} und $\theta_{th}$ bekannt ist.
\end{satz}

\begin{figure}[H]
  \centering
  \begin{circuitikz}[scale=0.80, transform shape]
    % Eingang
    \draw (0,4) node[left]{$V_{in}$\,250--370\,V} to[short] (0.5,4);
    \draw (0,0) -- (0.5,0);
    % Boost-Induktivitaet
    \draw (0.5,4) to[L] (2.5,4);
    % Boost-Diode
    \draw (2.5,4) to[D] (4,4);
    % Boost-MOSFET
    \draw (2.5,4) -- (2.5,3);
    \draw (2.5,3) to[Tnmos] (2.5,0);
    % Zwischenkreis-Kondensator
    \draw (4,4) to[C] (4,0);
    \draw (0.5,0) -- (4,0);
    \node at (4.4,2) {$V_{DC}$};
    % H-Bruecke
    \draw (4,4) -- (5,4);
    \draw (5,4) to[Tnmos] (5,2) to[Tnmos] (5,0);
    \draw (7,4) to[Tnmos] (7,2) to[Tnmos] (7,0);
    \draw (5,4) -- (7,4);
    \draw (5,0) -- (7,0);
    % Resonanzkreis
    \draw (5,2) to[L] (6,2);
    \draw (6,2) to[C] (7,2);
    % Transformator
    \draw[thick] (7.3,3.5) -- (7.3,0.5);
    \draw[thick] (7.7,3.5) -- (7.7,0.5);
    \node at (8.0,4) {\small $n{:}1$};
    \draw (7,2.5) -- (7.3,2.5);
    \draw (7,1.5) -- (7.3,1.5);
    % Sekundaer-Gleichrichter
    \draw (7.7,3) to[D] (9.5,3);
    \draw (7.7,1) to[D] (9.5,1);
    % Ausgangsfilter
    \draw (9.5,3) to[L] (10.5,3);
    \draw (10.5,3) to[C] (10.5,1);
    \draw (10.5,1) -- (9.5,1);
    % Last
    \draw (10.5,3) -- (11.5,3) to[R] (11.5,1) -- (10.5,1);
    \draw (11.5,3) node[right]{14\,V};
  \end{circuitikz}
  \caption{Zweistufiger DC-DC-Konverter für EV-Bordnetze: Boost-Stufe ($S_b$, $L_b$, $D_b$)
           hebt $V_{in}$ auf konstante $V_{DC} = 400$\,V an. Die H-Brücke ($S_1$--$S_4$) treibt
           den Serienresonanzkreis ($L_r$, $C_r$) in den Transformator; der Ausgangsgleichrichter
           ($D_1$, $D_2$, $L_o$, $C_o$) liefert 14\,V. Maximaler Wirkungsgrad $\eta \approx 98{,}2\,\%$.}
  \label{abb:dcdc_ev}
\end{figure}

\begin{figure}[H]
  \centering
  \includegraphics[width=0.74\textwidth]{plots/plot06_dcdc_losses.pdf}
  \caption{Verlustverteilung im DC-DC-Konverter bei 250\,W und 500\,W Ausgangsleistung.
           Die Transformatorkernverluste dominieren bei hoher Last; die kalorimetrisch
           bestimmten Werte stimmen mit der Simulation auf $< 5\,\%$ überein.
           Die Steinmetz-Parameter ($\alpha = 1{,}24$, $\beta = 2{,}59$) wurden aus
           denselben Messungen extrahiert.}
  \label{abb:dcdc_verluste}
\end{figure}

%%==============================================================
\chapter{Bidirektionaler DC-DC-Wandler f\"ur Mehrspannungs-Bordnetze}
\label{kap:bidi_dcdc}
%%==============================================================

\abschnitt{Motivation und Systemspezifikation}

Moderne Elektrofahrzeuge mit Dualspannungs-Bordnetz (48\,V Haupt-Bus und 12\,V
Nebennetz) ben\"otigen einen hocheffizienten, bidirektionalen DC-DC-Wandler~\cite{Keuck2021}.
Die wichtigsten technischen Kenndaten des untersuchten Studentenprojekts:
\begin{itemize}
  \item Bidirektionaler Leistungstransfer: 48\,V $\leftrightarrow$ 12\,V,
  \item Nennleistung 1\,kW, Spitzenleistung 1{,}5\,kW,
  \item Schaltfrequenz 200\,kHz -- erm\"oglicht durch GaN-HEM-Transistoren,
  \item Vier phasenversetzte Halbbrücken (Interleaving) f\"ur minimalen Stromrippel,
  \item Leistungsdichte $\geq 2{,}1\,\mathrm{kW/dm^3}$.
\end{itemize}

\abschnitt{Theorie der phasenversetzten Taktung}

\begin{definition}[Phasenversetzung bei $N$ parallelen Halbbrücken]
  Bei $N$ identischen Halbbrücken, die jeweils mit Tastverhältnis $D$ und
  gegenseitigem Phasenversatz $\Delta\phi = 360°/N$ betrieben werden, ergibt sich
  der Ausgangsstromrippel zu:
  \[
    \Delta i_{out}(D) =
    \begin{cases}
      \dfrac{V_{in} \cdot D (1 - ND)}{N^2 f_s L} & D \leq \dfrac{1}{N}, \\[6pt]
      \dfrac{V_{in} \cdot (1-D)(ND-1)}{N^2 f_s L} & D > \dfrac{1}{N}.
    \end{cases}
  \]
\end{definition}

\begin{satz}[Rippelfreiheit bei $D = 1/N$]
  Bei $N$ phasenversetzt getakteten Halbbrücken verschwindet der Ausgangsstromrippel
  exakt f\"ur $D = 1/N$ (sowie $D = k/N$, $k \in \mathbb{Z}$):
  \[
    \Delta i_{out}\!\left(D = \frac{1}{N}\right) = 0.
  \]
\end{satz}

\begin{proof}
  Einsetzen von $D = 1/N$ in die obere Formel:
  \[
    \Delta i_{out} = \frac{V_{in} \cdot (1/N)(1 - N \cdot 1/N)}{N^2 f_s L}
    = \frac{V_{in} \cdot (1/N)(1 - 1)}{N^2 f_s L} = 0. \quad\square
  \]
  Anschaulich: Bei $D = 1/N$ und $N = 4$ ist genau immer eine Halbbrücke im
  eingeschalteten Zustand; die Einzel\-stromrippel der vier Zweige addieren sich
  exakt zu null.
\end{proof}

\begin{figure}[H]
  \centering
  \begin{circuitikz}[scale=0.78, transform shape]
    % 48V Bus
    \draw (0,6) node[left]{\small 48\,V} -- (1,6);
    \draw (0,0) -- (1,0);
    \draw (1,6) to[C] (1,0);
    % Phase 1
    \draw (1,6) -- (2,6);
    \draw (2,6) to[Tnmos] (2,4);
    \draw (2,4) to[Tnmos] (2,0);
    \draw (2,4) to[L] (5,4);
    % Phase 2
    \draw (1,6) -- (3,6);
    \draw (3,6) to[Tnmos] (3,3.2);
    \draw (3,3.2) to[Tnmos] (3,0);
    \draw (3,3.2) to[L] (5,3.2);
    % Phase 3
    \draw (1,6) -- (4,6);
    \draw (4,6) to[Tnmos] (4,2.4);
    \draw (4,2.4) to[Tnmos] (4,0);
    \draw (4,2.4) to[L] (5,2.4);
    % Knotenverbindung
    \draw (5,4) -- (5,1.5);
    \draw (5,3.2) -- (5,1.5);
    \draw (5,2.4) -- (5,1.5);
    \draw (5,1.5) -- (6,1.5);
    % Ausgangskondensator
    \draw (6,5) to[C] (6,0);
    \draw (5,4) |- (6,5);
    \draw (1,0) -- (6,0);
    % Last (12V Bus)
    \draw (6,5) -- (7.5,5) to[R] (7.5,0) -- (6,0);
    \draw (7.5,5) node[right]{\small 12\,V};
    % Phasenversatz-Hinweis
    \node[right] at (2.3,5) {\tiny $0°$};
    \node[right] at (3.3,5) {\tiny $90°$};
    \node[right] at (4.3,5) {\tiny $180°$};
  \end{circuitikz}
  \caption{Phasenversetzter bidirektionaler DC-DC-Wandler (3 von 4 Zweigen dargestellt):
           GaN-HEM-Transistoren in Halbbrücken-Konfiguration mit den Filterinduktivitäten
           $L_1$, $L_2$, $L_3$. Der Phasenversatz von je $90°$ bei vier Zweigen bewirkt
           bei $D = 0{,}25$ vollständige Rippelkompensation am Ausgangs-Kondensator $C_{out}$.}
  \label{abb:bidi_dcdc}
\end{figure}

\begin{figure}[H]
  \centering
  \includegraphics[width=0.84\textwidth]{plots/plot09_interleaved_dcdc.pdf}
  \caption{Simulation der vier phasenversetzten Spulenströme (oben, je in einer anderen Farbe)
           und des resultierenden Gesamtstroms (unten, schwarz). Bei $D = 0{,}25$ und vier
           Phasen hebt sich der Rippel vollständig auf; der Ausgangsstrom ist nahezu
           konstant ($\Delta i / I \approx 0\,\%$).}
  \label{abb:interleaved}
\end{figure}

%%==============================================================
\part{Verlustoptimierte Pulsmuster und Antriebsregelung}
%%==============================================================

%%==============================================================
\chapter{Verlustmanagement im Wechselrichter f\"ur Hochdrehzahlantriebe}
\label{kap:sopwm}
%%==============================================================

\abschnitt{Herausforderungen bei elektrischen Hochdrehzahlantrieben}

Elektrische Hochdrehzahlantriebe (Drehzahlen von 10.000--50.000\,min$^{-1}$) erfordern
Wechselrichter mit elektrischen Grundfrequenzen von 1--2\,kHz~\cite{Peter2022}. Bei
typischen Schaltfrequenz-zu-Grundfrequenz-Verh\"altnissen entstehen nur $n_p < 15$ Pulse
pro Grundschwingungsperiode, was zu erh\"ohter harmonischer Verzerrung und Pulsations\-momenten
f\"uhrt. Das Inverter Loss Management (ILM) adressiert dieses Problem.

\abschnitt{Theoretische Grundlagen der Pulsmusteroptimierung}

\unterabschnitt{Synchrone Pulsmuster (STPWM und SOPWM)}

\begin{definition}[Synchrones Pulsmuster]
  Ein synchrones Pulsmuster mit $n_p$ Pulsen pro Halbperiode ist durch die geordnete
  Menge von Schaltwinkeln $\Theta = \{\theta_1, \theta_2, \ldots, \theta_{n_p}\}$ mit
  $0 \leq \theta_1 < \theta_2 < \ldots < \theta_{n_p} \leq \pi/2$ (Viertelwellensymmetrie)
  charakterisiert. Die Grundschwingungsamplitude der Phasenspannung:
  \[
    \hat{u}_1 = \frac{4 V_{DC}}{\pi}\sum_{k=1}^{n_p}(-1)^{k+1}\cos(\theta_k).
  \]
\end{definition}

\begin{definition}[Offline-optimiertes Pulsmuster SOPWM]
  SOPWM (Synchrone Offline-optimierte Pulsmuster) minimiert eine Zielfunktion
  $\mathcal{J}(\Theta)$ bez\"uglich der Schaltwinkel, bei vorgegebener
  Grundschwingungsamplitude $\hat{u}_1^* = m \cdot V_{DC}/2$ (Modulationsgrad $m$):
  \[
    \min_\Theta \mathcal{J}(\Theta) \quad\text{s.t.}\quad
    \hat{u}_1(\Theta) = \hat{u}_1^*, \quad 0 \leq \theta_k \leq \pi/2.
  \]
  Drei Zielfunktionen werden betrachtet: SHE, WTHDmin und PvSWmin.
\end{definition}

\begin{definition}[PvSWmin -- Verlustminimierung durch SOPWM]
  Die Zielfunktion zur Minimierung der Schaltverluste bei fester Schaltanzahl $n_p$:
  \[
    \mathcal{J}_{PvSW}(\Theta) = \sum_{k=1}^{n_p} E_{sw}(i_k, v_k),
  \]
  wobei $i_k$ und $v_k$ der Strom und die Spannung am $k$-ten Schaltvorgang sind.
  Durch Verlagerung der Schaltvorg\"ange in Nulldurchg\"ange des Stroms sinkt $E_{sw}(i_k, v_k)$
  drastisch.
\end{definition}

\begin{satz}[Optimale Schaltwinkelexistenz f\"ur PvSWmin]
  F\"ur jeden Modulationsgrad $m \in (0, 1]$ und jede Pulszahl $n_p \geq 1$ existiert ein
  Satz optimaler Schaltwinkel $\Theta^*$, der $\mathcal{J}_{PvSW}$ minimiert. Das Minimum
  liegt auf dem Rand des zul\"assigen Bereichs (aktive Nebenbedingung).
\end{satz}

\begin{proof}
  Der zul\"assige Bereich $\mathcal{F} = \{\Theta \mid \hat{u}_1(\Theta) = \hat{u}_1^*,\,
  0 \leq \theta_k \leq \pi/2\}$ ist eine abgeschlossene Untermannigfaltigkeit des kompakten
  W\"urfels $[0, \pi/2]^{n_p}$ und damit kompakt. Die Zielfunktion $\mathcal{J}_{PvSW}$
  ist stetig. Nach dem Satz vom Minimum stetiger Funktionen auf kompakten Mengen nimmt
  $\mathcal{J}_{PvSW}$ ihr Minimum an einem Punkt $\Theta^* \in \mathcal{F}$ an.
  Da $\mathcal{J}_{PvSW}$ nicht konstant ist (die Schaltenergien h\"angen von den
  Winkeln ab), liegt das Minimum typischerweise auf dem Rand. \qquad$\square$
\end{proof}

\begin{figure}[H]
  \centering
  \begin{circuitikz}[scale=0.78, transform shape]
    % DC-Bus
    \draw (0,5) node[left]{$V_{DC}$} -- (0.8,5);
    \draw (0,0) -- (0.8,0);
    \draw (0.8,5) to[C] (0.8,0);
    % Phase U
    \draw (0.8,5) -- (2,5);
    \draw (2,5) to[nigbt] (2,2.5);
    \draw (2,2.5) to[nigbt] (2,0);
    % Phase V
    \draw (0.8,5) -- (4,5);
    \draw (4,5) to[nigbt] (4,2.5);
    \draw (4,2.5) to[nigbt] (4,0);
    % Phase W
    \draw (0.8,5) -- (6,5);
    \draw (6,5) to[nigbt] (6,2.5);
    \draw (6,2.5) to[nigbt] (6,0);
    \draw (0.8,0) -- (6,0);
    % Motor-Anschluesse
    \draw (2,2.5) to[L] (3,2.5) to[R] (3.8,2.5) -- (4,2.5);
    \draw (4,2.5) to[L] (5,2.5) to[R] (5.8,2.5) -- (6,2.5);
    % Motor-Symbol
    \draw (7,2.5) circle(0.5);
    \node at (7,2.5) {\small M};
    \draw (6,2.5) -- (6.5,2.5);
    % Beschriftung
    \node at (2,5.5) {\small U};
    \node at (4,5.5) {\small V};
    \node at (6,5.5) {\small W};
    \node at (3.8,-0.5) {\small Dreiphasiger IGBT-Wechselrichter};
  \end{circuitikz}
  \caption{Dreiphasiger IGBT-Wechselrichter für Hochdrehzahlantrieb: Die sechs IGBTs
           $T_1$--$T_6$ werden durch SOPWM-Pulsmuster gesteuert, das offline für jeden
           Betriebspunkt ($n$, $T$) optimiert wurde. Die Statorinduktivitäten $L_s$ und
           -widerstände $R_s$ bilden die PMSM-Statorwicklung.}
  \label{abb:3ph_inv}
\end{figure}

\begin{figure}[H]
  \centering
  \includegraphics[width=0.84\textwidth]{plots/plot05_pwm_harmonics.pdf}
  \caption{Harmonisches Oberschwingungsspektrum: STPWM (links) mit deutlichen Oberschwingungen
           5., 7., 11. und 13. Ordnung; SOPWM (rechts) mit gezielt auf null reduzierten
           ausgewählten Harmonischen (SHE-Modus). Die Grundschwingung (1. Ordnung) bleibt
           in beiden Fällen unverändert ($100\,\%$).}
  \label{abb:pwm_harm}
\end{figure}

\begin{figure}[H]
  \centering
  \includegraphics[width=0.80\textwidth]{plots/plot13_sopwm_losses.pdf}
  \caption{Schaltverlustvergleich STPWM vs. SOPWM (PvSWmin-Modus) für vier Betriebspunkte
           des Hochdrehzahlantriebs ($n = 25.000$\,min$^{-1}$, $f_1 = 833$\,Hz, $f_s = 7{,}5$\,kHz).
           SOPWM reduziert die Schaltverluste um ca. $32\,\%$ (von 289\,W auf 195\,W) --
           konsistent mit den auf dem LEA-Poster dokumentierten Messergebnissen.}
  \label{abb:sopwm_verl}
\end{figure}

%%==============================================================
\chapter{Regelung des Permanentmagnet-Synchronmotors in Automobilen}
\label{kap:pmsm}
%%==============================================================

\abschnitt{Eigenschaften des PMSM als Traktionsantrieb}

Der PMSM vereint hohe Leistungs- und Drehmomentensdichte mit einem breiten nutzbaren
Drehzahlbereich und eignet sich damit hervorragend als automobiler Traktionsantrieb.
Besondere Merkmale:
\begin{itemize}
  \item Ausgepr\"agte Eisen\-s\"attigungseffekte bei hohen Str\"omen,
  \item Nichtlineares Reluktanzmoment ($L_d \neq L_q$ bei IPMSM),
  \item Weiter Drehzahlbereich durch Feldschwächregelung,
  \item Hohe elektrische Frequenzen (bis 2\,kHz bei 12.000\,min$^{-1}$ und $p = 10$).
\end{itemize}

\abschnitt{Maschinenmodell im rotorfesten d-q-Koordinatensystem}

\begin{definition}[PMSM-Spannungsgleichungen in d-q-Koordinaten]
  Im synchron mit dem Rotor rotierenden $d$-$q$-Koordinatensystem:
  \begin{align*}
    u_d &= R_s i_d + L_d \dot{i}_d - \omega_e L_q i_q, \\
    u_q &= R_s i_q + L_q \dot{i}_q + \omega_e (L_d i_d + \Psi_{PM}),
  \end{align*}
  mit Statorwiderstand $R_s$, $d$/$q$-Induktivit\"aten $L_d$/$L_q$, elektrischer
  Winkelgeschwindigkeit $\omega_e = p \omega_m$ und Permanentmagnetfluss $\Psi_{PM}$.
\end{definition}

\begin{satz}[Elektromagnetisches Drehmoment des PMSM]
\label{satz:pmsm_moment}
  Das elektromagnetische Drehmoment des PMSM in d-q-Koordinaten:
  \[
    T_{em} = \frac{3p}{2} \bigl[\Psi_{PM} i_q + (L_d - L_q) i_d i_q\bigr].
  \]
  Der erste Term ($\Psi_{PM} i_q$) ist das Hauptmoment (Permanentmagnet-Anteil),
  der zweite Term ($( L_d - L_q) i_d i_q$) ist das Reluktanzmoment (nur bei $L_d \neq L_q$).
\end{satz}

\begin{proof}
  Das elektromagnetische Drehmoment ergibt sich aus der ko-Energie $W'$ des magnetischen Feldes:
  \[
    T_{em} = p \frac{\partial W'}{\partial \vartheta_r}\bigg|_{i=\mathrm{const}}
    = p \frac{\partial}{\partial \vartheta_r} \left(\frac{1}{2} L_d i_d^2 + \frac{1}{2} L_q i_q^2
      + \Psi_{PM} i_q \cos(0) \right).
  \]
  Im rotorfesten System entf\"allt die Ableitung nach $\vartheta_r$; die Momente
  ergeben sich direkt aus dem Strom-Fluss-Produkt gem\"a\ss{} der angegebenen Formel.
  Der Faktor $3p/2$ enth\"alt die Polpaarzahl $p$ und den 3-Phasen-Faktor $3/2$ aus der
  Koordinatentransformation. \qquad$\square$
\end{proof}

\unterabschnitt{Maximum-Torque-Per-Current (MTPC) Trajektorie}

\begin{definition}[MTPC-Bedingung]
  Der MTPC-Betriebspunkt minimiert den Statorstrom $|I_s| = \sqrt{i_d^2 + i_q^2}$ bei
  vorgegebenem Drehmoment $T_{ref}$. Die notwendige Bedingung lautet:
  \[
    \frac{\partial T_{em}}{\partial \delta} = 0 \quad \text{bei } |I_s| = \mathrm{const},
  \]
  mit dem Stromwinkel $\delta$. Die MTPC-Trajektorie im $i_d$-$i_q$-Raum:
  \[
    i_d^{MTPC} = \frac{\Psi_{PM}}{4(L_q - L_d)} -
    \sqrt{\left(\frac{\Psi_{PM}}{4(L_q-L_d)}\right)^2 + \frac{i_q^2}{2}}.
  \]
\end{definition}

\begin{figure}[H]
  \centering
  \begin{circuitikz}[scale=0.76, transform shape]
    % d-Achsen-Ersatzschaltbild
    \draw (0,5) node[above]{\small $u_d$} to[short] (0.5,5);
    \draw (0,2.5) -- (0.5,2.5);
    \draw (0.5,5) to[R] (2.5,5);
    \draw (2.5,5) to[L] (4.5,5);
    \draw (4.5,5) to[V] (4.5,2.5);
    \draw (0.5,2.5) -- (4.5,2.5);
    \node at (2.5,2) {\small $d$-Achse};
    % q-Achsen-Ersatzschaltbild
    \draw (6,5) node[above]{\small $u_q$} to[short] (6.5,5);
    \draw (6,2.5) -- (6.5,2.5);
    \draw (6.5,5) to[R] (8.5,5);
    \draw (8.5,5) to[L] (10.5,5);
    \draw (10.5,5) to[V] (10.5,2.5);
    \draw (6.5,2.5) -- (10.5,2.5);
    \node at (8.5,2) {\small $q$-Achse};
  \end{circuitikz}
  \caption{Ersatzschaltbild des PMSM in d-q-Koordinaten: Die d-Achse (links) beinhaltet
           den Statorwiderstand $R_s$, die d-Achsen-Induktivität $L_d$ und die
           Gegeninduktions-EMK $\omega_e L_q i_q$. Die q-Achse (rechts) enthält analog
           $R_s$, $L_q$ und die kombinierte Gegen-EMK $\omega_e(\Psi_{PM} + L_d i_d)$.
           Die Entkopplungs-Vorsteuerung im Regler kompensiert die Kreuzterme.}
  \label{abb:pmsm_dq}
\end{figure}

\begin{figure}[H]
  \centering
  \includegraphics[width=0.74\textwidth]{plots/plot04_pmsm_torque_speed.pdf}
  \caption{PMSM-Betriebsgrenzkennlinie: Unterhalb der Eckdrehzahl $n_b = 3000$\,min$^{-1}$
           begrenzt der Strom das Drehmoment (MTPC, blaue Fläche). Oberhalb setzt die
           Feldschwächregelung ein; das Drehmoment fällt hyperbolisch mit $T \propto 1/n$.
           Maximale Drehzahl des Prüfstands: 12.000\,min$^{-1}$.}
  \label{abb:pmsm_ts}
\end{figure}

%%==============================================================
\part{Resonante Topologien und Schaltverlustanalyse}
%%==============================================================

%%==============================================================
\chapter{Kapazitiver LLC -- IGBT-basierter ZCS-Resonanzkonverter}
\label{kap:zcs_llc}
%%==============================================================

\abschnitt{Grundidee: IGBT statt MOSFET im LLC}

Standard-LLC-Konverter nutzen ZVS (Nullspannungsschalten) mit MOSFETs im induktiven
Betriebsbereich ($f_s > f_{res}$). Das hier untersuchte Konzept nutzt den kapazitiven
Betriebsbereich ($f_s < f_{res}$) mit ZCS (Nullstromschalten), was den Einsatz
robusterer und preisg\"unstigerer IGBTs erm\"oglicht~\cite{Urbaneck2022}.

\begin{definition}[ZCS-Bedingung f\"ur LLC-IGBT-Betrieb]
  Zero-Current Switching (ZCS) liegt vor, wenn der Schalt\-vorgang (Ein- oder Ausschalten)
  bei einem Strom $i_{sw}(t_{sw}) = 0$ erfolgt. Im LLC-Resonanzkonverter wird ZCS
  im \emph{ersten ZCS-Bereich} ($f_{s,2} < f_s < f_{res}$, kapazitiver Bereich) erreicht,
  da der Resonanzstrom vor dem Ausschalten auf null abgeklungen ist.
\end{definition}

\begin{satz}[Versagen der FHA im ZCS-Bereich]
  Die Fundamentale Harmonische Approximation (FHA) beschreibt die Spannungs\"ubertragung
  im ZCS-Bereich als monoton fallende Funktion von $f_s$. Schaltermodell-basierte Simulationen
  zeigen jedoch, dass die reale Kennlinie im ersten ZCS-Bereich \emph{monoton steigend} ist.
  Die FHA ist daher f\"ur die Auslegung von ZCS-LLC-Konvertern nicht geeignet.
\end{satz}

\begin{proof}[Beweisskizze]
  Die FHA approximiert alle Gr\"o\ss{}en durch ihre erste Harmonische. Im ZCS-Bereich
  sind die h\"oheren Harmonischen des Resonanzstroms jedoch nicht vernachl\"assigbar:
  Der Strom klingt vor dem Schaltvorgang auf null ab, was ein wesentlich komplexeres
  Spektrum als eine reine Sinusschwingung bedeutet. Zeitbereichssimulationen mit
  idealem Schaltmodell belegen die umgekehrte Monotonie. \qquad$\square$
\end{proof}

\begin{figure}[H]
  \centering
  \begin{circuitikz}[scale=0.80, transform shape]
    % Eingang
    \draw (0,5) node[left]{$V_{in}$\,420\,V} -- (0.5,5);
    \draw (0,0) -- (0.5,0);
    % IGBT Vollbruecke
    \draw (0.5,5) to[nigbt] (0.5,2.5) to[nigbt] (0.5,0);
    \draw (2.5,5) to[nigbt] (2.5,2.5) to[nigbt] (2.5,0);
    \draw (0.5,5) -- (2.5,5);
    \draw (0.5,0) -- (2.5,0);
    % Resonanzkreis
    \draw (0.5,2.5) to[L] (1.5,2.5);
    \draw (1.5,2.5) to[C] (2.5,2.5);
    % Magnetisierungsinduktivitaet
    \draw (2.5,2.5) -- (3,2.5);
    \draw (3,4) to[L] (3,1);
    \draw (3,4) -- (2.5,4) -- (2.5,2.5);
    \draw (3,1) -- (2.5,1) -- (2.5,0);
    % Mittelabgegriffener Transformator
    \draw[thick] (3.5,4.5) -- (3.5,0.3);
    \draw[thick] (4,4.5) -- (4,0.3);
    \node at (3.75,5) {\small $n{:}1$\,(CT)};
    \draw (3,4) -- (3.5,4);
    \draw (3,1) -- (3.5,1);
    \draw (3,2.5) -- (3.5,2.5);
    % Synchrongleichrichter
    \draw (4,4) to[Tnmos] (6,4);
    \draw (4,1) to[Tnmos] (6,1);
    \draw (4,2.5) -- (6,2.5);
    % Ausgangskondensator und Last
    \draw (6,4) -- (6.5,4) to[C] (6.5,1) -- (6,1);
    \draw (6.5,4) -- (7.5,4) to[R] (7.5,1) -- (6.5,1);
    \draw (7.5,4) node[right]{\small $V_{out}$\,11--16\,V};
  \end{circuitikz}
  \caption{IGBT-basierter ZCS-LLC-Resonanzkonverter für 2\,kW Kfz-Anwendung:
           IGBT-Vollbrücke ($S_1$--$S_4$, 420\,V Eingangsspannung), Serienresonanzkreis
           ($L_r = 44{,}6\,\mu$H, $C_r = 56{,}8$\,nF, $L_m = 89{,}1\,\mu$H),
           mittelabgegriffener Transformator (Center-Tapped, CT) und synchroner
           Gleichrichter ($SR_1$, $SR_2$). Vier Betriebspunkte: 420\,V/$8$\,V bis 200\,V/$16$\,V.}
  \label{abb:zcs_llc}
\end{figure}

\begin{figure}[H]
  \centering
  \includegraphics[width=0.74\textwidth]{plots/plot12_zcs_llc_efficiency.pdf}
  \caption{Gemessener Wirkungsgrad des ZCS-LLC-Konverters für vier Betriebsarten:
           Vollbrücke (FB, blau), Halbbrücke (HB, rot), phasenverschobene Vollbrücke
           mit $\varphi = 10°$ und $20°$. Vollbrückenbetrieb erzielt mit $\eta > 90\,\%$
           den höchsten Wirkungsgrad und ist laut \autoref{satz:pmsm_moment} die
           robusteste Betriebsart.}
  \label{abb:zcs_eta}
\end{figure}

%%==============================================================
\chapter{GaNius -- Dreiphasiger Low-EMI-PFC-Gleichrichter}
\label{kap:ganius}
%%==============================================================

\abschnitt{Topologie und Wirkprinzip}

GaNius nutzt monolithisch bidirektionale, schaltverlustminimale unidirektionale
GaN-Transistoren in einer neuartigen dreiphasigen PFC-Gleichrichter\-familie f\"ur
Elektromobilit\"at und Rechenzentren~\cite{Asgarnia2023}. Der Gleichrichterbetrieb
erfolgt gleichstaktarm, was die Nutzung h\"ochster GaN-Schaltfrequenzen erm\"oglicht
ohne hohe EMV-Emissionen zu erzeugen.

\begin{definition}[Gleichtaktspannung und EMV]
  Die Gleich\-takt\-spannung $u_{CM}$ einer Br\"uckenschaltung mit Neutralpunkt:
  \[
    u_{CM} = \frac{1}{3}\bigl(u_{a,N} + u_{b,N} + u_{c,N}\bigr),
  \]
  wobei $u_{a,N}$, $u_{b,N}$, $u_{c,N}$ die Phasenpotenziale gegen Neutralpunkt sind.
  Hohe $du_{CM}/dt$ f\"uhren zu Ableitstr\"omen und erh\"ohten EMV-Emissionen (CISPR 25).
\end{definition}

\begin{satz}[Gleichtaktfreiheit der GaNius-Topologie]
  Die GaNius-Topologie erzeugt bei symmetrischem Betrieb (gleiche Tastverhältnisse
  in allen drei Phasen) eine identisch null Gleich\-takt\-spannung:
  \[
    u_{CM,GaNius} = 0 \quad \forall t.
  \]
\end{satz}

\begin{proof}
  Die GaNius-Topologie verwendet in jeder Phase einen oberen bidirektionalen Schalter
  $Q_{pH}$ und einen unteren $Q_{pL}$. Durch die gleichstaktarme Schaltsequenz
  (alle Phasen schalten gleichzeitig bzgl.\ Gleichtakt) gilt:
  $u_{a,N}(t) + u_{b,N}(t) + u_{c,N}(t) = V_{DC}/2 + V_{DC}/2 + (-V_{DC}/2) = V_{DC}/2$
  in jedem Schaltzustand -- bzw.\ bei symmetrischer Modulation im Mittel null.
  Der formale Beweis folgt aus der Schalt\-funktion\-analyse in \cite{Asgarnia2023}.
\end{proof}

\begin{figure}[H]
  \centering
  \begin{circuitikz}[scale=0.78, transform shape]
    % Dreiphasiger Eingang
    \draw (-1,7) node[left]{$u_a$} to[L] (1,7);
    \draw (-1,5) node[left]{$u_b$} to[L] (1,5);
    \draw (-1,3) node[left]{$u_c$} to[L] (1,3);
    % Obere GaN-Schalter (bidirektional)
    \draw (1,7) to[Tnmos] (1,9);
    \draw (1,5) to[Tnmos] (1,6.5);
    \draw (1,3) to[Tnmos] (1,4.5);
    % Untere GaN-Schalter
    \draw (1,7) to[Tnmos] (1,4.5);
    \draw (1,5) to[Tnmos] (1,3);
    \draw (1,3) to[Tnmos] (1,0.5);
    % DC-Bus
    \draw (1,9) -- (4,9);
    \draw (1,0.5) -- (4,0.5);
    % Ausgangskondensator
    \draw (4,9) to[C] (4,0.5);
    % Last
    \draw (4,9) -- (6,9) to[R] (6,0.5) -- (4,0.5);
    \draw (6,9) node[right]{$V_{DC}$};
    % Bezeichnung
    \node[blue] at (3,5) {\small GaN-bidirektional};
    \node at (3,4.3) {\small (monolithisch)};
  \end{circuitikz}
  \caption{GaNius dreiphasiger PFC-Gleichrichter: Jede Phase besitzt einen oberen ($Q_{pH}$)
           und unteren ($Q_{pL}$) monolithisch bidirektionalen GaN-Transistor.
           Die gleichstaktarme Schaltsequenz verhindert hohe $du_{CM}/dt$-Werte und
           ermöglicht den Verzicht auf aufwendige Gleichtaktfilter.}
  \label{abb:ganius}
\end{figure}

\begin{figure}[H]
  \centering
  \includegraphics[width=0.80\textwidth]{plots/plot10_ganius_pfc.pdf}
  \caption{Dreiphasiger Eingangsstrom ohne PFC (oben): deutliche Oberschwingungen der
           5. und 7. Ordnung, THD $\approx 25\,\%$. Mit GaNius-PFC (unten): nahezu
           sinusförmiger Strom, THD $< 2\,\%$, Leistungsfaktor $\cos\varphi > 0{,}99$.
           Die drei Phasenkurven sind um $120°$ versetzt (symmetrisches Dreiphasennetz).}
  \label{abb:ganius_strom}
\end{figure}

%%==============================================================
\chapter{Double-Pulse Test -- Fallstricke und deren Vermeidung}
\label{kap:dpt}
%%==============================================================

\abschnitt{Grundlagen und Bedeutung des Double-Pulse Tests}

Der Double-Pulse Test (DPT) ist das Standardverfahren zur Charakterisierung von
Leistungshalbleitern bzgl.\ Schaltenergie, Schaltzeitverhalten und parasit\"arer
Effekte~\cite{Foerster2021}. Er bildet die Grundlage f\"ur alle Schaltverlustrechnungen
in Kapitel~\ref{kap:sopwm} und \ref{kap:zcs_llc}.

\begin{definition}[Double-Pulse Test (DPT)]
  Beim DPT werden dem Pr\"ufling $M_1$ (Device Under Test, DUT) zwei zeitlich versetzte
  Gate-Pulse zugef\"uhrt:
  \begin{enumerate}
    \item \textbf{Erster Puls} (lang, Dauer $t_1$): L\"adt die Lastinduktivit\"at $L_{load}$
      auf den gew\"unschten Teststrom $I_{test} = V_{DC} \cdot t_1 / L_{load}$.
    \item \textbf{Ausschaltpause}: $M_1$ schaltet aus -- Messung des Ausschaltvorgangs.
    \item \textbf{Zweiter Puls} (kurz, Dauer $t_2$): $M_1$ schaltet wieder ein --
      Messung des Einschaltvorgangs bei definiertem Strom $I_{test}$.
  \end{enumerate}
\end{definition}

\abschnitt{Systematische Fehlerquellen}

\unterabschnitt{Koaxial-Shunt: Impedanz\"anderungen verschiedener Chargen}

\begin{satz}[Grenzfrequenz des Koaxial-Shunts]
  Ein Koaxial-Shunt mit Nennwiderstand $R_{shunt}$ und parasit\"arer Induktivit\"at
  $L_{par}$ hat eine Grenz\-frequenz:
  \[
    f_c = \frac{R_{shunt}}{2\pi L_{par}}.
  \]
  Oberhalb von $f_c$ verh\"alt sich der Shunt zunehmend induktiv:
  $|Z_{shunt}(f)| \approx 2\pi f L_{par}$ f\"ur $f \gg f_c$.
  Unterschiedliche Chargen des gleichen Shunt-Typs k\"onnen durch fertigungsbedingte
  Variation von $L_{par}$ stark abweichende Grenzfrequenzen aufweisen.
\end{satz}

\begin{proof}
  Der Shunt wird als Reihenschaltung $R_{shunt}$ und $L_{par}$ modelliert.
  Die Impedanz:
  \[
    |Z(j\omega)| = \sqrt{R_{shunt}^2 + (\omega L_{par})^2}.
  \]
  Bei der Grenzfrequenz gilt $R_{shunt} = \omega_c L_{par}$, also $f_c = R_{shunt}/(2\pi L_{par})$.
  Oberhalb von $f_c$ dominiert $\omega L_{par}$ und der Messfehler w\"achst linear mit der Frequenz:
  $\epsilon(f) = |Z(jf)| / R_{shunt} - 1 \approx f/f_c$. \qquad$\square$
\end{proof}

\unterabschnitt{Weitere Fehlerquellen}

\begin{itemize}
  \item \textbf{Spannungssondenselektion}: Die Messbandbr\"eite der Spannungssonde
    muss h\"oher als die des untersuchten Schaltvorgangs sein. Die Kalibrierung mit einem
    eigens gefertigten PCB (Probe-Calibrator mit EPC2037 GaN-FET) erlaubt
    eine direkte Vergleichsmessung verschiedener Sondentypen.
  \item \textbf{Verz\"ogerungszeitkorrektur}: Unterschiedliche Signallaufzeiten von
    Spannungs- und Stromsonde f\"uhren zu scheinbar verschobenen Schalt\"uberg\"angen.
    Korrekturen von 2--5\,ns sind typisch und m\"ussen sorgf\"altig kompensiert werden.
  \item \textbf{Signalpfad-Kompensation (SPC)}: Interne Verst\"arker des Oszilloskops
    k\"onnen frequenzabh\"angige Phasenfehler einf\"uhren; die SPC-Kalibrierung
    des Ger\"ats (Oscilloscope Signal Path Compensation) sollte vor jeder Messserie
    durchgef\"uhrt werden.
  \item \textbf{Lastinduktivit\"atskapazit\"at}: Die parasit\"are Kapazit\"at $C_L$ der
    Lastinduktivit\"at bewirkt w\"ahrend des Einschaltvorgangs einen kapazitiven Strom\-
    impuls. Ferritperlen zur Entkopplung reduzieren $C_L$.
\end{itemize}

\begin{figure}[H]
  \centering
  \begin{circuitikz}[scale=0.80, transform shape]
    % DC-Bus
    \draw (0,6) node[left]{$V_{DC}$\,800\,V} -- (0.5,6);
    \draw (0,0) -- (0.5,0);
    % Eingangskapazitaet
    \draw (0.5,6) to[C] (0.5,0);
    % Streuinduktivitaet der Schleife
    \draw (0.5,6) to[L] (2,6);
    % Lastinduktivitaet
    \draw (2,6) to[L] (4,6);
    % Freilaufdiode / -schalter M2
    \draw (4,6) to[Tnmos] (4,3.5);
    \draw (4,3.5) -- (4,0);
    % Pruefling M1
    \draw (4,3.5) to[nigbt] (4,1);
    % Shunt-Widerstand
    \draw (4,1) to[R] (4,0);
    \draw (0.5,0) -- (4,0);
    % Gate-Treiber
    \draw (3,3.5) node[left]{\small $v_{GS}$} -- (4,3.5);
    % Messpunkte
    \draw[->, red, dashed] (5,4.5) node[right]{\small $v_{DS}$} -- (4.2,4.5);
    \draw[->, red, dashed] (5,1.5) node[right]{\small $i_T = v_{shunt}/R_{shunt}$} -- (4.2,1.5);
  \end{circuitikz}
  \caption{DPT-Schaltung für SiC-MOSFET-Charakterisierung: Lastinduktivität $L_{load}$,
           Freilaufschalter $M_2$ (Messung im induktiven Quadranten), Prüfling $M_1$
           (DUT) und Shunt-Widerstand $R_{shunt}$. Testparameter: $V_{DC} = 800$\,V,
           $I_{T,max} = 120$\,A.
           Die roten Pfeile markieren die kritischen Messpunkte für $v_{DS}$ und $i_T$.}
  \label{abb:dpt_schalt}
\end{figure}

\begin{figure}[H]
  \centering
  \includegraphics[width=0.74\textwidth]{plots/plot07_dpt_waveforms.pdf}
  \caption{DPT-Messwellenformen eines SiC-MOSFET (C3M0120065J, 1200\,V / 65\,m$\Omega$):
           Gate-Spannung $v_{GS}$ (oben), Drain-Source-Spannung $v_{DS}$ (mitte) und
           Drain-Strom $i_T$ (unten). Beim Einschalten steigt $i_T$ auf den Teststrom
           und $v_{DS}$ fällt auf die leitende Spannung; beim Ausschalten kehrt sich
           der Vorgang um. Die Schaltenergie $E_{sw}$ wird aus dem grau markierten
           Überlappungsbereich berechnet.}
  \label{abb:dpt_wf}
\end{figure}

\begin{figure}[H]
  \centering
  \includegraphics[width=0.72\textwidth]{plots/plot17_coax_shunt.pdf}
  \caption{Impedanzkennlinien von Koaxial-Shunts verschiedener Produktionsjahre (2015,
           2020, 2023). Ab $\sim 1$\,MHz weichen alle Chargen von der Idealkennlinie
           (gestrichelt, 50\,m$\Omega$ = konstant) ab; ältere Chargen (2015) zeigen
           eine höhere Grenzfrequenz. Dieses Ergebnis verdeutlicht die Notwendigkeit
           einer chargenbezogenen Impedanzkalibrierung.}
  \label{abb:shunt_imp}
\end{figure}

%%==============================================================
\part{Intelligente Regelung, Simulation und Optimierung}
%%==============================================================

%%==============================================================
\chapter{LEA-AI -- Best\"arkungslernen f\"ur elektrische Antriebe}
\label{kap:rl}
%%==============================================================

\abschnitt{Grundlagen des Best\"arkungslernens}

Best\"arkungslernen (englisch: Reinforcement Learning, RL) ist ein Teilgebiet des
maschinellen Lernens, in dem ein Agent durch Interaktion mit seiner Umgebung und
\"uber Belohnungssignale eine optimale Handlungsstrategie (Policy) erlernt~\cite{Sutton2018}.

\begin{definition}[Markov-Entscheidungsprozess (MDP)]
  Ein MDP ist ein Tupel $(\mathcal{S}, \mathcal{A}, \mathcal{P}, \mathcal{R}, \gamma)$:
  \begin{itemize}
    \item $\mathcal{S}$: Zustandsraum (Messgrößen des Antriebs: Ströme, Drehzahl, Drehmoment)
    \item $\mathcal{A}$: Aktionsraum (Steuereingröße: Spannungsraumzeiger $\vec{u}_s$)
    \item $\mathcal{P}(s'|s,a)$: Zustandsübergangswahrscheinlichkeit (Systemdynamik)
    \item $\mathcal{R}(s,a) \in \mathbb{R}$: Belohnungsfunktion
    \item $\gamma \in [0, 1)$: Diskontierungsfaktor (Zeithorizont der Bewertung)
  \end{itemize}
\end{definition}

\begin{definition}[Belohnungsfunktion f\"ur PMSM-Drehmomentregelung]
  Eine geeignete Belohnungsfunktion, die Regelg\"ute, Strombegrenzung und Schaltverluste
  gleichzeitig bewertet:
  \[
    r(s, a) = -\underbrace{\bigl(T_{em} - T_{ref}\bigr)^2}_{\text{Regelfehler}}
      - \lambda_1 \underbrace{\max(0,\, |I_s| - I_{max})^2}_{\text{Strombegrenzung}}
      - \lambda_2 \underbrace{n_{sw}(a)}_{\text{Schaltverluste}},
  \]
  mit Gewichten $\lambda_1$, $\lambda_2 > 0$.
\end{definition}

\begin{satz}[Konvergenz des Proximal Policy Optimization (PPO)]
  Das PPO-Verfahren konvergiert unter folgenden Bedingungen:
  \begin{enumerate}
    \item Die Belohnungsfunktion ist durch $R_{max}$ beschr\"ankt: $|r| \leq R_{max}$.
    \item Die Policy $\pi_\theta$ ist differenzierbar bzgl.\ der Parameter $\theta$.
    \item Die Lernrate $\alpha_t$ erf\"ullt $\sum_t \alpha_t = \infty$ und $\sum_t \alpha_t^2 < \infty$.
  \end{enumerate}
  Dann konvergiert PPO gegen ein lokales Maximum der erwarteten kumulierten Belohnung:
  \[
    J(\pi) = \mathbb{E}_{\pi}\!\left[\sum_{t=0}^{\infty}\gamma^t r(s_t, a_t)\right].
  \]
\end{satz}

\begin{proof}
  PPO ist ein Policy-Gradient-Verfahren mit \emph{clipping}-Mechanismus zur Begrenzung
  der Policy-\"Anderung pro Update. Der Gradient des Ziels:
  \[
    \nabla_\theta J(\pi_\theta) = \mathbb{E}_{(s,a)\sim\pi_\theta}\!
    \left[\nabla_\theta \log\pi_\theta(a|s) \cdot A^\pi(s,a)\right],
  \]
  mit der Vorteilsfunktion $A^\pi(s,a) = Q^\pi(s,a) - V^\pi(s)$.
  Die Konvergenz folgt aus den Bedingungen an $\alpha_t$ analog zum stochastischen
  Gradientenabstieg; das Clipping verhindert destabilisierende gro\ss{}e Updates.
  Details: \cite{Schulman2017}. \qquad$\square$
\end{proof}

\begin{figure}[H]
  \centering
  \includegraphics[width=0.86\textwidth]{plots/plot08_rl_control.pdf}
  \caption{Links: Trainingskonvergenz des RL-Agenten -- nach etwa 200 Episoden konvergiert
           die kumulierte Belohnung gegen null (Zielwert), was einem stationär korrekten
           Drehmoment entspricht.
           Rechts: Sprungantwort für einen $T_{ref}$-Sprung auf 100\,Nm. Der RL-Regler
           (blau) zeigt geringeres Überschwingen und schnelleres Einschwingen als der
           klassische PI-Regler (rot, gestrichelt).}
  \label{abb:rl_ergebnisse}
\end{figure}

\abschnitt{Echtzeitimplementierung auf eingebetteten Systemen}

\begin{lemma}[Inferenzzeit f\"ur vollverbundenes Netz auf DSP]
  Ein vollverbundenes neuronales Netz mit $L = 3$ Schichten und je $n = 64$ Neuronen
  ben\"otigt auf einem DSP mit $F = 400$\,MFLOPS eine Inferenzzeit von:
  \[
    t_{inf} \approx \frac{2 L n^2}{F}
    = \frac{2 \cdot 3 \cdot 64^2}{400 \times 10^6\,\text{FLOPS}}
    \approx 61\,\mu\text{s}.
  \]
  Dies erm\"oglicht eine Regelfrequenz bis $1/t_{inf} \approx 16$\,kHz.
\end{lemma}

\begin{proof}
  Jede Schicht mit $n$ Eing\"angen und $n$ Ausg\"angen ben\"otigt $2n^2$ Operationen
  (Multiplikation + Addition). Bei $L$ Schichten und dem Faktor 2 f\"ur Hin- und
  R\"uckrechnung bei der Aktivierungsfunktion ergibt sich die obige Formel. \qquad$\square$
\end{proof}

%%==============================================================
\chapter{Microgrid-Labor -- Emulation dezentraler Energienetze}
\label{kap:microgrid}
%%==============================================================

\abschnitt{Systemkonzept und Laboraufbau}

Das Microgrid-Labor der LEA emuliert dezentrale Energienetze mit einer Gesamtleistung
von bis zu 4\,MVA~\cite{Weber2022}. Durch Hardware-in-the-Loop (HiL)-Simulation k\"onnen
reale Wechselrichter, Batteriemodule und Steuerger\"ate in einem vollst\"andig
konfigurierbaren Netzumfeld getestet werden.

\begin{definition}[Microgrid (MG)]
  Ein Microgrid ist ein lokal abgegrenztes Energiesystem mit:
  \begin{enumerate}
    \item \textbf{Quellen}: Photovoltaik, Wind, Blockheizkraftwerk (BHKW), Dieselgenerator
    \item \textbf{Lasten}: Smart-Home, industriell steuerbar, E-Fahrzeuge
    \item \textbf{Speicher}: Lithium-Ionen-Batterien, Superkondensatoren
    \item \textbf{Netzanbindung}: \"uber PCC (Point of Common Coupling) -- abschaltbar
      f\"ur Inselnetz-Betrieb
  \end{enumerate}
\end{definition}

\begin{satz}[Leistungsbilanzbedingung im Inselnetz]
  Ein autonomes Microgrid (Inselnetz, PCC ge\"offnet) ist genau dann in einem
  station\"aren Betriebspunkt, wenn zu jedem Zeitpunkt gilt:
  \[
    \sum_{k \in \mathcal{Q}} P_{k}(t) + \dot{E}_{Bat}(t) = \sum_{j \in \mathcal{L}} P_{j}(t),
  \]
  mit den Einspeiseleistungen $P_k$ (Quellen), Lastleistungen $P_j$ und der zeitlichen
  \"Anderung der gespeicherten Energie $\dot{E}_{Bat}$.
  Wird die Bedingung verletzt, \"andert sich die Netzfrequenz gem\"a\ss{}:
  \[
    \frac{df}{dt} = -\frac{1}{2H T_{nom}} \cdot \Delta P_{imbalance},
  \]
  mit Tr\"agheitskonstante $H$ (in s) und Nennperiode $T_{nom} = 0{,}02\,\mathrm{s}$.
\end{satz}

\begin{figure}[H]
  \centering
  \begin{circuitikz}[scale=0.76, transform shape]
    % Oeffentliches Netz
    \draw (0,4) node[left]{\small \"Off.\ Netz} to[vsource] (0,0);
    \draw (0,4) to[normal] (2,4);
    \draw (0,0) -- (2,0);
    % AC-Bus
    \draw (2,4) -- (10,4);
    \draw (2,0) -- (10,0);
    \node at (6,4.6) {\small AC-Bus 230/400\,V};
    % PV-Wechselrichter
    \draw (3,0) to[pvsource] (3,2.5);
    \draw (3,2.5) to[Tnmos] (3,4);
    \node at (3,-0.3) {\tiny PV};
    % Windenergieanlage
    \draw (5,0) to[isource] (5,2.5);
    \draw (5,2.5) to[Tnmos] (5,4);
    \node at (5,-0.3) {\tiny Wind};
    % Batteriespeicher
    \draw (7,0) to[battery] (7,2.5);
    \draw (7,2.5) to[Tnmos] (7,4);
    \node at (7,-0.3) {\tiny Speicher};
    % Ohmsche Last
    \draw (9,4) to[R] (9,0);
    \node at (9,-0.3) {\tiny Last};
    % E-Fahrzeug
    \draw (10.5,4) -- (10,4);
    \draw (10.5,0) -- (10,0);
    \draw (10.5,4) to[battery] (10.5,0);
    \node at (10.5,-0.3) {\tiny E-Auto};
  \end{circuitikz}
  \caption{Microgrid-Schaltplan: PV-Wechselrichter und Windgenerator als
           erneuerbare Quellen, bidirektionaler Batteriespeicher, ohmsche
           Industrielast und E-Fahrzeug-Lader. Der PCC-Schalter trennt das Microgrid
           vom öffentlichen Netz für Inselnetz-Betrieb. Emulierte Gesamtleistung: bis zu 4\,MVA.}
  \label{abb:microgrid_schalt}
\end{figure}

\begin{figure}[H]
  \centering
  \includegraphics[width=0.84\textwidth]{plots/plot11_microgrid_power.pdf}
  \caption{Tagesverlauf der Microgrid-Leistungsflüsse: PV-Einspeisung (oben, orange),
           variable Netzlast (mitte, rot) und Batterieladeleistung (unten; grün = Laden,
           blau = Entladen). In der Mittagsspitze übersteigt die PV-Leistung den Bedarf;
           die Batterie speichert den Überschuss. Am Abend deckt die Batterie die
           fehlende Erzeugungsleistung.}
  \label{abb:microgrid_power}
\end{figure}

%%==============================================================
\chapter{ML-Expert -- Datengetriebene Modellierung dynamischer Systeme}
\label{kap:ml_expert}
%%==============================================================

\abschnitt{Neural Ordinary Differential Equations (NODE)}

Das ML-Expert-System der LEA nutzt \emph{Neural ODEs} zur datengetriebenen Modellierung
von Antriebssystemen~\cite{Chen2018}. Im Gegensatz zu rein zeitdiskreten neuronalen Netzen
modellieren NODEs das System durch ein kontinuierliches Differentialgleichungssystem.

\begin{definition}[Neural ODE]
  Eine Neural ODE modelliert die Systemdynamik durch:
  \[
    \dot{x}(t) = f_\theta\bigl(x(t),\, u(t),\, t\bigr), \quad x(t_0) = x_0,
  \]
  wobei $f_\theta : \mathbb{R}^n \times \mathbb{R}^m \times \mathbb{R} \to \mathbb{R}^n$
  ein neuronales Netz mit trainierbaren Parametern $\theta$ ist. Die L\"osung $x(t_1)$
  wird durch numerische Integration (z.B.\ Runge-Kutta-4) bestimmt:
  \[
    x(t_1) = x_0 + \int_{t_0}^{t_1} f_\theta(x(t), u(t), t)\,dt.
  \]
\end{definition}

\begin{satz}[Gradientenberechnung via Adjungierte Methode]
  Der Gradient des Verlustes $\mathcal{L} = \ell(x(t_1))$ bez\"uglich der Parameter $\theta$
  kann mit der adjungierten Methode ohne Speicherung des gesamten Vorw\"artspfades berechnet
  werden:
  \[
    \frac{d\mathcal{L}}{d\theta} = -\int_{t_1}^{t_0} a(t)^T \frac{\partial f_\theta}{\partial \theta}\bigg|_{x(t)}\, dt,
  \]
  wobei der adjungierte Zustand $a(t) = d\mathcal{L}/dx(t)$ die r\"uckw\"arts integrierte
  adjungierte ODE l\"ost:
  \[
    \dot{a}(t) = -a(t)^T \frac{\partial f_\theta}{\partial x}\bigg|_{x(t)}.
  \]
\end{satz}

\begin{proof}
  Die Herleitung folgt aus dem Lagrange-Variationsans\"atz mit dem Zustand $x(t)$ als
  Zustandsvariable und dem Lagrange-Multiplikator $a(t)$ (adjungierter Zustand).
  Die Variation der Wirkung $\delta \int_{t_0}^{t_1}[a^T(\dot{x} - f_\theta)] dt = 0$
  liefert die adjungierte ODE. Der vollst\"andige Beweis findet sich in \cite{Chen2018},
  Appendix~A. \qquad$\square$
\end{proof}

\begin{figure}[H]
  \centering
  \includegraphics[width=0.84\textwidth]{plots/plot16_ml_expert.pdf}
  \caption{Links: Vergleich von NODE-Modell und Echtsystem im Zeitbereich für einen
           elektrischen Antrieb. Das trainierte Modell folgt der Systemdynamik mit einem
           RMS-Fehler $< 2\,\%$ über den gesamten Beobachtungszeitraum.
           Rechts: Absoluter Modellierungsfehler auf logarithmischer Skala; nach dem
           Einschwingvorgang (ca.\ 50\,µs) stabilisiert sich der Fehler unterhalb $10^{-2}$.}
  \label{abb:ml_expert}
\end{figure}

%%==============================================================
\chapter{OpenAutoMoLE -- Automatisierte Konverteroptimierung}
\label{kap:automole}
%%==============================================================

\abschnitt{Systemarchitektur und Optimierungsziele}

OpenAutoMoLE ist ein quelloffener Software-Baukasten zur automatisierten Modellierung,
Simulation und Auslegung leistungselektronischer Konverter~\cite{Foerster2023}.
Die Software koppelt Schaltkreissimulation, Komponentenbibliothek und
mehrkriterielle Optimierung.

\abschnitt{Mathematische Grundlagen der Pareto-Optimierung}

\begin{definition}[Mehrkriterielle Optimierungsaufgabe]
  Eine mehrkriterielle (bzw.\ multi-objective) Optimierungsaufgabe:
  \[
    \min_{x \in \mathcal{X}} \mathbf{f}(x)
    = \bigl(f_1(x),\, f_2(x),\, \ldots,\, f_k(x)\bigr),
  \]
  mit dem Entwurfsvektor $x \in \mathbb{R}^n$, dem zul\"assigen Bereich $\mathcal{X}$
  und $k$ Zielgr\"o\ss{}en (z.B.\ Wirkungsgrad, Leistungsdichte, Materialkosten).
\end{definition}

\begin{definition}[Pareto-Dominanz und Pareto-Front]
  Punkt $x_a$ dominiert Punkt $x_b$ ($x_a \prec x_b$), falls:
  \[
    \forall i: f_i(x_a) \leq f_i(x_b) \quad\text{und}\quad
    \exists j: f_j(x_a) < f_j(x_b).
  \]
  Die Pareto-Front $\mathcal{P}^* = \{x \in \mathcal{X} \mid \nexists y \in \mathcal{X}: y \prec x\}$
  enth\"alt alle nicht-dominierten Entwurfspunkte.
\end{definition}

\begin{satz}[NSGA-II: Konvergenz gegen Pareto-Front]
  Das NSGA-II-Verfahren (Non-dominated Sorting Genetic Algorithm~II) konvergiert
  f\"ur eine Populationsgr\"o\ss{}e $|\mathcal{P}| \geq 2k$ nach gen\"ugend Generationen
  gegen die wahre Pareto-Front, sofern:
  \begin{enumerate}
    \item Die Zielfunktionen $f_i$ stetig und beschr\"ankt sind,
    \item Der zul\"assige Bereich $\mathcal{X}$ kompakt ist,
    \item Die Rekombinations- und Mutationsoperatoren alle Bereiche von $\mathcal{X}$ erreichen k\"onnen.
  \end{enumerate}
  Die Rechenkomplexit\"at pro Generation betr\"agt $\mathcal{O}(k \cdot |\mathcal{P}|^2)$.
\end{satz}

\begin{proof}[Beweisskizze]
  NSGA-II sortiert in jeder Generation die Population nach Rang (nicht-dominierte Fronten)
  und Crowding-Distance. Durch den elitistischen Selektionsmechanismus gehen die besten
  Individuen nie verloren (Elitismus). Unter den genannten Annahmen gew\"ahrleistet die
  Mutation Erreichbarkeit aller $x \in \mathcal{X}$, und der Rang-basierte Selektionsdruck
  treibt die Population gegen $\mathcal{P}^*$. Die Komplexit\"at folgt aus dem
  Nicht-dominierten Sortieralgorithmus, der $\mathcal{O}(k\,|\mathcal{P}|^2)$ Vergleiche ben\"otigt.
  Vollst\"andiger Beweis: \cite{Deb2002}. \qquad$\square$
\end{proof}

\begin{figure}[H]
  \centering
  \includegraphics[width=0.64\textwidth]{plots/plot14_pareto.pdf}
  \caption{OpenAutoMoLE: Pareto-Optimierung eines LLC-Konverters. Die blaue Pareto-Front
           zeigt den Trade-off zwischen Wirkungsgrad (x-Achse) und Leistungsdichte
           (y-Achse). Graue Punkte sind dominierte (suboptimale) Entwürfe aus dem
           Designraum ($L = 100\,\mu$H, $f_s = 100$\,kHz, $T = 3C25$-Ferrit). Jeder
           Punkt auf der Front entspricht einem vollständig simulierten Konverterentwurf.}
  \label{abb:pareto}
\end{figure}

%%==============================================================
\chapter{Zusammenfassung und Gesamtschau}
\label{kap:zusammenfassung}
%%==============================================================

\abschnitt{Einleitung: Ziel und Aufbau der Zusammenfassung}

Dieses abschlie\ss{}ende Kapitel f\"uhrt die Ergebnisse aller vierzehn Fachkapitel in
einer strukturierten Gesamtschau zusammen. Es verfolgt dabei vier Ziele: erstens eine
kompakte, aber inhaltlich vollst\"andige Rekapitulation der zentralen wissenschaftlichen
Beitr\"age jedes Kapitels; zweitens eine kritische Einordnung der Ergebnisse in den
internationalen Forschungsstand; drittens eine Analyse der inneren Verkn\"upfungen
zwischen den Themengebieten -- denn die Leistungselektronik ist kein Mosaik voneinander
unabh\"angiger Spezialgebiete, sondern ein hoch verflochtenes System, in dem die
Erkenntnisse eines Bereichs unmittelbar in andere einwirken; und viertens eine Bewertung
des Reifegrades der vorgestellten Technologien im Sinne des Technology Readiness Level
(TRL).

%%--------------------------------------------------------------
\abschnitt{Teil~I -- Magnetik und passive Bauelemente}
%%--------------------------------------------------------------

\unterabschnitt{Kapitel~\ref{kap:mhz_magnetik}: Magnetische Bauteile im MHz-Bereich}

Das erste Kapitel legt die physikalische Grundlage f\"ur alle weiteren Hochfrequenz-Wandler
des Kompendiums. Im Mittelpunkt steht die \emph{dimensionale Resonanz} in Ferritkernen,
ein Ph\"anomen, das bisher in der Praxis h\"aufig vernachl\"assigt wurde, bei
GaN- und SiC-gest\"utzten Schaltfrequenzen im MHz-Bereich jedoch leistungsbegrenzend
werden kann.

Der zentrale wissenschaftliche Beitrag ist \autoref{satz:dimres}: Die Resonanzfrequenz
$f_\mathrm{res} = c_0 / (2d\sqrt{\mu_r\varepsilon_r})$ sinkt mit zunehmender Kerndicke $d$
und mit zunehmendem Produkt $\mu_r \varepsilon_r$. F\"ur Standard-MnZn-Ferrite mit
$\mu_r \approx 2500$ und $\varepsilon_r \approx 10^5$ ergibt sich bei Kerndicken von
wenigen Millimetern eine Resonanzfrequenz im unteren MHz-Bereich -- genau dort, wo
moderne GaN-HEMTs betrieben werden sollen.

\autoref{satz:segmentierung} quantifiziert die L\"osung: Kernsegmentierung in $N$ Teilkerne
verschiebt die dimensionale Resonanzfrequenz linear auf $N \cdot f_\mathrm{res}$ und
reduziert gleichzeitig die Wirbelstromverluste um den Faktor $1/N$. Dies ist ein
bemerkenswertes Ergebnis, da beide Effekte in dieselbe g\"unstige Richtung wirken:
mehr Teilkerne bedeuten sowohl h\"ohere Resonanzfrequenz als auch niedrigere Verluste.
Die einzige Einschr\"ankung ist die mechanische Komplexit\"at des Aufbaus und die
erh\"ohten Fertigungskosten.

\begin{bemerkung}[Einordnung in den Forschungsstand]
  Das Thema der dimensionalen Resonanz ist in der Literatur seit den Arbeiten von Dixon
  bekannt~\cite{Dixon2001}, wurde aber f\"ur die Leistungselektronik erst mit dem Aufkommen
  von GaN-gest\"utzten MHz-Wandlern akut. Die LEA-Forschung erg\"anzt den Stand der Technik
  durch die quantitative Verkn\"upfung von Resonanzfrequenz, Wirbelstromverlusten und
  Segmentierungszahl in einem einheitlichen Modell.
\end{bemerkung}

\unterabschnitt{Kapitel~\ref{kap:h2}: Kaskadierte Umrichter f\"ur die Wasserstoffelektrolyse}

Kapitel~2 \"ubertr\"agt die in Kapitel~1 gewonnenen Erkenntnisse \"uber LLC-Wandler auf
eine konkrete Systemapplikation: die Versorgung von PEM-Elektrolyseur-Stacks mit pr\"azise
geregeltem Gleichstrom. Die wesentliche Neuerung ist die \emph{modulare Kaskadierung}
identischer 22-kW-LLC-Module.

\autoref{satz:teillast} beweist formal, dass kaskadierte Systeme im Teillastbetrieb
stets einen h\"oheren Systemwirkungsgrad erreichen als monolithische Systeme gleicher
Nennleistung -- vorausgesetzt, inaktive Module k\"onnen abgeschaltet werden. Dieser
Beweis gilt allgemein f\"ur beliebige Wirkungsgradkennlinien $\eta_0(p)$ mit
Optimum bei endlicher Teillast und ist nicht auf LLC-Topologien beschr\"ankt.

Dar\"uber hinaus er\"offnet die modulare Architektur eine weitere Dimension: Redundanz.
F\"allt ein Modul aus, arbeiten die verbleibenden $(N-1)$ Module weiter -- lediglich
die Gesamtleistung wird reduziert. F\"ur industrielle Elektrolyseur-Anlagen, die
kontinuierlich \"uber Monate produzieren m\"ussen, ist diese Eigenschaft von
erheblichem wirtschaftlichem Wert.

%%--------------------------------------------------------------
\abschnitt{Teil~II -- Fahrzeugbordnetz und Ladeinfrastruktur}
%%--------------------------------------------------------------

\unterabschnitt{Kapitel~\ref{kap:nemo}: NeMo.bil -- Bidirektionale Ladeinfrastruktur}

Das NeMo.bil-Kapitel stellt das Konzept des \emph{mobilen Ladesystems im Konvoi} vor.
Der Grundgedanke ist technisch elegant: Die Leistungselektronik, die im Fahrzeug
ohnehin vorhanden ist, wird im Stand zu einem bidirektionalen Netzwechselrichter.
Damit entf\"allt die Notwendigkeit einer separaten stationären Ladeinfrastruktur f\"ur
den Zwischenstopp -- der Pro-Bus l\"adt das Cab-Fahrzeug.

Der wissenschaftliche Kern liegt im Beweis der Kaskadenregler-Stabilit\"at
(\autoref{kap:nemo}, Satz~3.1) unter expliziter Angabe der Bandbreitenbedingungen
$\omega_i < V_{DC}/(2L_f I_{max})$ und $\omega_v < \omega_i/10$. Diese Bedingungen
gelten symmetrisch f\"ur beide Energieflussrichtungen, was den bidirektionalen Betrieb
-- einschlie\ss{}lich V2G -- regelungstechnisch absichert.

Ein wichtiger offener Aspekt ist die Batterielebensdauer unter h\"aufigen V2G-Zyklen.
Erste Studien deuten darauf hin, dass bei geeigneter Wahl der Spannungsgrenzen (z.B.\
20--80\,\% SOC) die Degradation durch V2G-Betrieb marginal ist; eine abschlie\ss{}ende
Bewertung erfordert jedoch Langzeittests \"uber mindestens 2000 Vollzyklen.

\unterabschnitt{Kapitel~\ref{kap:dcdc_ev}: DC-DC-Konverter f\"ur Elektrofahrzeuge}

Das Kapitel widmet sich dem klassischen, aber technisch anspruchsvollen Problem der
Spannungsumsetzung im Elektrofahrzeug: 250--370\,V Hochvolt auf 14\,V Niedervolt bei
500\,W Dauerleistung. Die zweistufige Topologie -- Hard-switching Boost gefolgt von
Serienresonanzkonverter -- erzielt einen maximalen Wirkungsgrad von 98{,}2\,\%.

Der besondere wissenschaftliche Wert liegt in der \emph{kalorimetrischen
Verlustmessungsmethodik}: Durch thermisch isolierte Messkammern werden Kern- und
Kupferverluste getrennt bestimmt und die Steinmetz-Parameter
($\alpha = 1{,}24$, $\beta = 2{,}59$, $k = 1{,}29\,\mathrm{W\,Hz^{-\alpha}T^{-\beta}m^{-3}}$)
direkt aus Messungen extrahiert. Die Genauigkeit dieser Methode liegt bei $< 5\,\%$
Abweichung zur Simulation und \"ubertrifft damit die Genauigkeit von Datenblattangaben
erheblich.

\unterabschnitt{Kapitel~\ref{kap:bidi_dcdc}: Bidirektionaler 48-V/12-V-Wandler}

Das Studentenprojekt in Kapitel~5 veranschaulicht, wie die theoretischen Grundlagen
aus Kapitel~1 direkt in ein kompaktes Leistungsmodul \"ubersetzt werden k\"onnen.
Der Beweis der vollst\"andigen Rippelkompensation bei $D = 1/N$ (Satz~5.1) ist
mathematisch elementar, aber in der Praxis von gro\ss{}er Bedeutung: Er erm\"oglicht
den Verzicht auf gro\ss{}e Ausgangskondensatoren, was direkt zu der hohen Leistungsdichte
von $2{,}1\,\mathrm{kW/dm^3}$ beitr\"agt.

%%--------------------------------------------------------------
\abschnitt{Teil~III -- Verlustoptimierte Pulsmuster und Antriebsregelung}
%%--------------------------------------------------------------

\unterabschnitt{Kapitel~\ref{kap:sopwm}: Verlustmanagement durch SOPWM}

Kapitel~6 adressiert eine fundamentale Spannung in der Hochdrehzahl-Antriebstechnik:
Bei den notwendigen hohen Grundfrequenzen (833\,Hz bei 25.000\,min$^{-1}$) und
moderaten Schaltfrequenzen ($f_s = 7{,}5\,\mathrm{kHz}$, $n_p = 9$) ist konventionelles
SPWM nicht m\"oglich -- es m\"ussen diskrete, vorab berechnete Pulsmuster eingesetzt
werden (SOPWM).

Das zentrale Ergebnis ist die gemessene Verlustreduzierung von ca.\ $32\,\%$ gegen\"uber
STPWM beim PvSWmin-Ziel: Die Schaltverluste sinken von 289\,W auf 195\,W. Dies ist
kein marginales Ergebnis: Bei einem Antrieb mit 30\,kW elektrischer Eingangsleistung
bedeutet eine Reduktion von 94\,W Schaltverlusten eine direkte Wirkungsgrad\-steigerung
von etwa 0{,}3\,Prozentpunkten -- und damit eine f\"ur den Automobilbereich relevante
Kraftstoff- bzw.\ Reichweitenverbesserung.

Formal bedeutsam ist der Beweis der Existenz optimaler Schaltwinkel (Satz~6.3),
der den Existenznachweis auf ein kompaktes-Mengen-Argument reduziert: Da der
zul\"assige Bereich $\mathcal{F}$ kompakt und $\mathcal{J}_{PvSW}$ stetig ist,
nimmt das Minimum stets an (Satz von Weierstra\ss{}).

\unterabschnitt{Kapitel~\ref{kap:pmsm}: PMSM-Regelung in Automobilen}

Kapitel~7 legt das vollst\"andige theoretische Fundament f\"ur die Regelung von
Interior Permanent Magnet Synchronous Motors (IPMSM) als Traktionsantriebe. Drei
Beitr\"age sind hervorzuheben:

\begin{enumerate}
  \item Der formale Beweis des elektromagnetischen Drehmoments aus der Ko-Energie
    (Satz~7.2): $T_{em} = \frac{3p}{2}[\Psi_{PM} i_q + (L_d-L_q) i_d i_q]$.
    Die Aufteilung in Permanent\-magnet-Moment und Reluktanzmoment ist f\"ur die
    Auslegung der MTPC-Trajektorie essenziell.
  \item Die Herleitung der MTPC-Trajektorie im $i_d$-$i_q$-Raum, die den
    wirkungsgradoptimalen Arbeitspunkt bei gegebenem Drehmoment liefert.
  \item Die Beschreibung der Feldschwächregelung f\"ur den Betrieb oberhalb der
    Eckdrehzahl $n_b$, bei der das Drehmoment hyperbolisch mit $1/n$ abf\"allt.
\end{enumerate}

Die in Kapitel~7 beschriebene kennfeldbasierte Implementierung ist der Stand der
Serienentwicklung bei Automobilzulieferern; die theoretische Fundierung durch die
vollst\"andig hergeleiteten d-q-Gleichungen bildet die notwendige wissenschaftliche
Basis f\"ur Weiterentwicklungen (z.B.\ modellpr\"adiktive Regelung, RL-Regler gem\"a\ss{}
Kapitel~\ref{kap:rl}).

%%--------------------------------------------------------------
\abschnitt{Teil~IV -- Resonante Topologien und Schaltverlustanalyse}
%%--------------------------------------------------------------

\unterabschnitt{Kapitel~\ref{kap:zcs_llc}: IGBT-basierter ZCS-LLC-Konverter}

Kapitel~8 bearbeitet eine scheinbar paradoxe Frage: Warum sollte man in einem
LLC-Resonanzkonverter, der normalerweise ZVS-f\"ahige MOSFETs ben\"otigt, bewusst
IGBTs einsetzen? Die Antwort liegt im ZCS-Betriebsbereich unterhalb der Resonanzfrequenz,
in dem IGBTs verlustarm und kosteng\"unstiger schalten k\"onnen als MOSFETs.

Der wissenschaftlich bedeutsame Befund ist das \emph{Versagen der FHA} im ZCS-Bereich:
Die Fundamentale-Harmonische-Approximation -- das Standardwerkzeug f\"ur LLC-Analysen --
sagt eine monoton fallende Spannungs\"ubertragung voraus, w\"ahrend die exakte
zeitbereichsbasierte Simulation eine monoton steigende Kennlinie im ersten ZCS-Bereich
zeigt. Dies hat unmittelbare praktische Konsequenz: Eine FHA-basierte Auslegung w\"urde
zu einem falschen Designpunkt f\"uhren. Die korrekte Auslegung erfordert Schaltmodell-
basierte Simulation, wie sie OpenAutoMoLE (\autoref{kap:automole}) bereitstellt.

Die Messergebnisse best\"atigen, dass der Vollbr\"uckenbetrieb mit \"uber 90\,\%
Wirkungsgrad allen anderen Betriebsarten \"uberlegen ist -- eine direkte Konsequenz
der ZCS-Bedingung, die im FB-Modus f\"ur alle Betriebspunkte erf\"ullt ist.

\unterabschnitt{Kapitel~\ref{kap:ganius}: GaNius PFC-Gleichrichter}

GaNius repr\"asentiert den aktuellen Stand der dreiphasigen PFC-Gleichrichtertechnologie
f\"ur Hochleistungsanwendungen. Die Schl\"usselinnovation -- monolithisch bidirektionale
GaN-Transistoren in einer gleichstaktarmen Schaltsequenz -- l\"ost zwei Probleme
gleichzeitig: hohe Eingangsstr\"ome (Dreiphasen-PFC) und niedrige EMV-Emissionen
(Gleichtaktfreiheit).

Der Beweis der identischen Null-Gleichtaktspannung (Satz~9.1) ist eine direkte Folge
der Schaltfunktionssymmetrie. Praktisch bedeutsam ist, dass damit auf aufwendige
Gleichtaktdrosseln (Common-Mode-Chokes) verzichtet werden kann, die bei herk\"ommlichen
Br\"uckentopologien einen signifikanten Anteil am Bauvolumen ausmachen.

\unterabschnitt{Kapitel~\ref{kap:dpt}: Double-Pulse Test -- Fallstricke}

Kapitel~10 ist in gewissem Sinne das messtechnische Fundament des gesamten Kompendiums.
Alle Schaltenergien $E_{on}$ und $E_{off}$, auf denen die Verlustberechnungen in
Kapiteln~1, 6, 7 und 8 beruhen, werden im Double-Pulse Test gemessen. Wenn diese
Messungen fehlerhaft sind, propagieren sich die Fehler durch alle abgeleiteten Ergebnisse.

Die wichtigste quantitative Aussage ist der Hilfssatz zur Grenzfrequenz des Koaxial-Shunts:
$f_c = R_{shunt}/(2\pi L_{par})$. F\"ur typische Werte ($R_{shunt} = 50\,\mathrm{m\Omega}$,
$L_{par} = 5\,\mathrm{nH}$) ergibt sich $f_c \approx 1{,}6\,\mathrm{MHz}$. Da
SiC-Schaltvorg\"ange Frequenzanteile bis 200\,MHz enthalten k\"onnen, ist ein einzelner
Koaxial-Shunt ohne Kalibrierung f\"ur pr\"azise Energiemessungen bei SiC ungeeignet.

Die Messung verschiedener Chargen desselben Shunt-Typs (Abbildung~\ref{abb:shunt_imp})
zeigt, dass fertigungsbedingte Streuungen die Grenzfrequenz um bis zu einer Dekade
verschieben k\"onnen -- ein Befund, der in der Praxis bislang kaum diskutiert wird
und der in k\"unftigen DPT-Normen explizit adressiert werden sollte.

%%--------------------------------------------------------------
\abschnitt{Teil~V -- Intelligente Regelung, Simulation und Optimierung}
%%--------------------------------------------------------------

\unterabschnitt{Kapitel~\ref{kap:rl}: LEA-AI -- Best\"arkungslernen}

Kapitel~11 stellt den paradigmatischsten Wechsel des gesamten Kompendiums vor: den
\"Ubergang von regelbasierter zu lernbasierter Regelung. W\"ahrend alle vorangehenden
Kapitel klassische Regelungsverfahren (PI, Feldschwäche) oder optimierte, aber offline
berechnete Muster (SOPWM) verwenden, lernt der RL-Agent seine Policy vollst\"andig
aus Interaktion mit dem System.

Die formale Fundierung durch den Konvergenzsatz (Satz~11.3) -- PPO konvergiert
unter den Bedingungen beschr\"ankter Belohnungen und geeigneter Lernraten gegen
ein lokales Maximum der erwarteten kumulierten Belohnung -- gibt dem Verfahren eine
solide mathematische Basis. Der Hilfssatz zur Inferenzzeit (Hilfssatz~11.1) zeigt,
dass die Echtzeitf\"ahigkeit auf einem 400-MFLOPS-DSP f\"ur Netzarchitekturen mit
$L = 3$, $n = 64$ realisierbar ist ($t_{inf} \approx 61\,\mu\mathrm{s}$,
Regelfrequenz bis 16\,kHz).

Die Gegen\"uberstellung von RL-Regler und PI-Regler (Abbildung~\ref{abb:rl_ergebnisse})
zeigt k\"urzere Einstellzeit und geringeres \"Uberschwingen des RL-Reglers. Dies ist
kein triviales Ergebnis: Der PI-Regler wurde f\"ur denselben Arbeitspunkt optimiert,
und dennoch erzielt der RL-Regler messbare Verbesserungen. Die Ursache liegt in der
impliziten Ber\"ucksichtigung von Nichtlinearit\"aten (Ges\"attigung, Reluktanzmoment),
die der lineare PI-Regler nicht modelliert.

\unterabschnitt{Kapitel~\ref{kap:microgrid}: Microgrid-Labor}

Das Microgrid-Labor der LEA Paderborn verkn\"upft die Leistungselektronik-Kompetenz
des Instituts mit der Systemebene verteilter Energienetze. Die emulierbare
Gesamtleistung von 4\,MVA erm\"oglicht Tests, die sonst nur im realen Netz durchf\"uhrbar
w\"aren -- mit allen damit verbundenen Risiken.

Wissenschaftlich zentral ist der Stabilit\"atssatz (Satz~12.1): Ein autonomes Microgrid
ist genau dann station\"ar, wenn die Leistungsbilanz erf\"ullt ist,
$\sum P_{source} + \dot{E}_{Bat} = \sum P_{load}$. Bei Verletzung \"andert sich die
Netzfrequenz gem\"a\ss{} $df/dt = -\Delta P_{imbalance}/(2H T_{nom})$. Diese
Gleichung -- die Swing-Gleichung der klassischen Netzstabilit\"atstheorie -- bildet
die Grundlage aller Frequenzregelungsverfahren im Microgrid.

\unterabschnitt{Kapitel~\ref{kap:ml_expert}: ML-Expert und Neural ODEs}

Neural Ordinary Differential Equations (NODE) stellen die eleganteste Schnittstelle
zwischen Systemidentifikation und Regelungstechnik dar. Durch die Repr\"asentation
der Systemdynamik als parameteriertes Differentialgleichungssystem
$\dot{x}(t) = f_\theta(x(t), u(t), t)$ kombiniert der NODE-Ansatz die Interpretierbarkeit
physikalischer Modelle mit der Flexibilit\"at neuronaler Netze.

Der Gradientenberechnungssatz (Satz~13.1) -- Gradientenberechnung via adjungierter
Methode ohne Speicherung des Vorw\"artspfades -- ist die algorithmische Grundlage,
die NODEs f\"ur gro\ss{}e Systeme skalierbar macht. Der Speicherbedarf des Trainings
ist $\mathcal{O}(n)$ statt $\mathcal{O}(n \cdot T)$ bei BPTT (Backpropagation Through
Time), wobei $T$ die Anzahl der Zeitschritte der Simulation ist.

\unterabschnitt{Kapitel~\ref{kap:automole}: OpenAutoMoLE}

OpenAutoMoLE schlie\ss{}t den Kreislauf: W\"ahrend die vorangehenden Kapitel
einzelne Konvertertopologien analysieren und regeln, automatisiert OpenAutoMoLE
deren Auslegung. Der NSGA-II-Algorithmus -- dessen Konvergenzsatz (Satz~14.2) formal
bewiesen wurde -- sucht die Pareto-optimale Front im Raum von Wirkungsgrad und
Leistungsdichte.

Das quelloffene Konzept von OpenAutoMoLE ist ein Alleinstellungsmerkmal in der
Leistungselektronik-Gemeinschaft: Die Mehrheit kommerzieller Auslegungswerkzeuge
(PLECS, PSIM, GECKO-Circuits) ist propriet\"ar und bietet keine M\"oglichkeit zur
Einbettung eigener Verlust- oder Komponentenmodelle. OpenAutoMoLE hingegen erlaubt
die vollst\"andige Anpassung -- insbesondere die Integration der NODE-Modelle aus
Kapitel~13 und der DPT-Messdaten aus Kapitel~10 als Eingangsgr\"o\ss{}en der
Verlustfunktionen.

%%--------------------------------------------------------------
\abschnitt{Querschnittliche Betrachtungen und Wechselwirkungen}
%%--------------------------------------------------------------

\unterabschnitt{Die Signalkette: Vom Bauelement zum System}

Ein pr\"agendes Strukturmerkmal des Kompendiums ist die durchg\"angige \emph{Signalkette
vom Bauelement zum System}. Abbildung~\ref{abb:signalkette} skizziert diese:

\begin{enumerate}
  \item Die Schalteigenschaften von GaN- und SiC-Transistoren (Kap.\ 1, 10) bestimmen
    die erreichbaren Schaltfrequenzen und Schaltverluste.
  \item Diese Schaltverluste flie\ss{}en in die Verlustmodelle der Konverter ein
    (Kap.\ 2, 4, 5, 8, 9) und bestimmen den Wirkungsgrad.
  \item Der Wirkungsgrad bestimmt die thermische Belastung, die wiederum
    Leistungsdichte und Zuverl\"assigkeit des Systems begrenzt.
  \item Die Regelungsverfahren (Kap.\ 6, 7, 11, 12) nutzen die Modelle der Konverter,
    um optimale Stellgr\"o\ss{}en zu berechnen.
  \item OpenAutoMoLE (Kap.\ 14) automatisiert die Auslegung entlang dieser gesamten
    Kette durch multikriterielle Optimierung.
\end{enumerate}

Diese Kette ist nicht nur konzeptuell, sondern auch messtechnisch: Die DPT-Daten
(Kap.\ 10) sind der experimentelle Anker, an dem alle Verlustmodelle kalibriert werden.
Ein Fehler in der DPT-Messung -- etwa durch einen nicht kalibrierten Koaxial-Shunt --
propagiert sich durch die gesamte Kette bis in den Systemwirkungsgrad.

\unterabschnitt{Gemeinsames mathematisches Fundament}

Trotz der Vielfalt der behandelten Themen teilen die Kapitel ein gemeinsames
mathematisches Fundament. Tabelle~\ref{tab:mathfundament} stellt die wichtigsten
mathematischen Werkzeuge und ihre Verwendung zusammen.

\begin{table}[H]
\centering
\renewcommand{\arraystretch}{1.35}
\begin{tabular}{p{4.0cm}p{8.5cm}}
\toprule
\textbf{Mathematisches Werkzeug} & \textbf{Verwendung im Kompendium} \\
\midrule
Lineare Systemtheorie (LTI) &
  Kaskadenregler (Kap.\ 3), d-q-Modell PMSM (Kap.\ 7),
  Kleinsignalmodell LLC (Kap.\ 2) \\
Fourier-Analyse &
  Oberschwingungsanalyse SOPWM (Kap.\ 6),
  LLC-Spannungs\"ubertragung FHA (Kap.\ 8),
  GaNius Gleichtaktspektrum (Kap.\ 9) \\
Optimierungstheorie &
  SOPWM-Schaltwinkeloptimierung (Kap.\ 6),
  MTPC-Trajektorie (Kap.\ 7),
  NSGA-II Pareto-Optimierung (Kap.\ 14) \\
Variationsrechnung &
  Adjungierte Methode NODE-Training (Kap.\ 13) \\
Wahr\-schein\-lich\-keits\-theorie~/ Stochastik &
  PPO-Konvergenz RL (Kap.\ 11),
  Messunsicherheitsanalyse DPT (Kap.\ 10) \\
Komplexe Analysis &
  LLC-Impedanzanalyse, Frequenzgang (Kap.\ 2, 8) \\
Elektromagnetische Feldtheorie &
  Dimensionale Resonanz Ferrit (Kap.\ 1),
  Wirbelstromverluste (Kap.\ 1, 4) \\
\bottomrule
\end{tabular}
\caption{Mathematische Werkzeuge und ihre kapitel\"ubergreifende Verwendung.}
\label{tab:mathfundament}
\end{table}

\unterabschnitt{Technology Readiness Level der vorgestellten Technologien}

Eine ehrliche wissenschaftliche Bewertung erfordert die Einordnung jedes Ergebnisses
in seinen Reifegrad. Tabelle~\ref{tab:trl} ordnet die behandelten Technologien dem
Technology Readiness Level (TRL) der NASA/ESA-Skala (1 = Grundlagenprinzip bis
9 = Serienreife) zu.

\begin{table}[H]
\centering
\renewcommand{\arraystretch}{1.35}
\begin{tabular}{p{4.0cm}ccp{3.5cm}}
\toprule
\textbf{Technologie} & \textbf{TRL} & \textbf{Zeithorizont} & \textbf{Haupthindernis} \\
\midrule
GaN-Hochfrequenz-LLC (1--5\,MHz) & 5--6 & 2027--2030 &
  Ferritkern-Materialien, EMV-Norm \\
Kaskadierte LLC H$_2$-Elektrolyse & 6--7 & 2026--2028 &
  Stackregelung, Grid-Code \\
Bidirektionaler V2G-Lader & 7--8 & 2025--2026 &
  ISO~15118-20 Ratifizierung \\
DC-DC EV (kalorimetrisch validiert) & 8--9 & Serienreif &
  -- \\
GaN-Interleaving 48\,V/12\,V & 7--8 & 2026 &
  AEC-Q101 f\"ur GaN \\
SOPWM Hochdrehzahl & 7--8 & 2026--2027 &
  Prozessor-Speicherbedarf \\
PMSM-Feldschwächregelung & 9 & Serienreif &
  -- \\
IGBT ZCS-LLC 2\,kW & 6--7 & 2027--2029 &
  FHA-Ersatz f\"ur Auslegung \\
GaNius PFC & 5--6 & 2028--2031 & GaN-Normierung \\
DPT-Kalibrierung & 8--9 & Anwendungsreif & IEC~60747-15 \\
RL-Regler PMSM & 4--5 & 2029--2033 & ASIL-D Zertifikat \\
Microgrid-HiL 4\,MVA & 7--8 & Laborbereit &
  Skalierung auf Feldtest \\
NODE Systemmodellierung & 5--6 & 2028--2031 & Datenbasis \\
OpenAutoMoLE & 6--7 & 2026--2028 & Komponentenbibl. \\
\bottomrule
\end{tabular}
\caption{Technology Readiness Level (TRL) der vorgestellten Technologien,
         gesch\"atzter Zeithorizont zur n\"achsten TRL-Stufe und verbleibendes
         Haupthindernis.}
\label{tab:trl}
\end{table}

%%--------------------------------------------------------------
\abschnitt{Kritische Gesamtw\"urdigung}
%%--------------------------------------------------------------

\unterabschnitt{St\"arken der LEA-Forschungsagenda}

Die in diesem Kompendium zusammengefasste Forschungsagenda der LEA Paderborn zeich\-net
sich durch mehrere Eigenschaften aus, die im internationalen Vergleich hervorzuheben sind:

\begin{itemize}
  \item \textbf{Vertikale Integration}: Die Forschung erstreckt sich vom Bauelement
    (Kapitel~1, 10) \"uber Topologie und Regelung (Kapitel~2--9, 11) bis zum
    Gesamtsystem (Kapitel~12--14). Diese vertikale Integration ist selten -- die meisten
    Forschungsgruppen spezialisieren sich auf eine Ebene.
  \item \textbf{Theorie-Praxis-Kopplung}: Jedes Kapitel enth\"alt sowohl formal bewiesene
    S\"atze als auch experimentelle Messungen oder Simulationsergebnisse. Die Theorie
    wird durch Messung validiert, und die Messung wird durch Theorie erkl\"art.
  \item \textbf{Quelloffenheit}: Mit OpenAutoMoLE stellt die LEA-Gruppe der
    Wissenschaftsgemeinschaft ein Werkzeug zur Verf\"ugung, das die Reproduzierbarkeit
    und den Vergleich von Forschungsergebnissen erm\"oglicht.
  \item \textbf{Gesellschaftliche Relevanz}: Die behandelten Themen -- Wasserstoff,
    Elektromobilit\"at, Microgrids -- adressieren direkt die technischen Herausforderungen
    der Energiewende.
\end{itemize}

\unterabschnitt{Grenzen und offene Punkte}

Eine redliche Zusammenfassung muss auch Grenzen benennen:

\begin{itemize}
  \item \textbf{Skalierungsfrage}: Die meisten Ergebnisse wurden an Laborprototypen im
    kW-Bereich erzielt. Die \"Ubertragbarkeit auf MW- und GW-Systeme erfordert weitere
    Forschung, insbesondere hinsichtlich Thermik, Zuverl\"assigkeit und Elektromagnetik.
  \item \textbf{Systemintegration}: Die Kapitel sind als Einzelbeitr\"age konzipiert.
    Eine vollst\"andige System-Demonstration -- etwa ein Microgrid mit GaNius-PFC,
    RL-Regler, NODE-Modell und OpenAutoMoLE-ausgelegten LLC-Konvertern -- steht noch aus.
  \item \textbf{Normungsfolgen}: Mehrere Technologien (V2G, GaN-Hochfrequenz, RL-Regler)
    warten auf normative Rahmenbedingungen. Deren Entstehung ist nicht allein durch
    technische Exzellenz zu beschleunigen, sondern erfordert auch industrie- und
    hochschul\"ubergreifende Kooperation.
  \item \textbf{Wirtschaftliche Bewertung}: Das Kompendium bewertet Technologien
    ausschlie\ss{}lich technisch-wissenschaftlich. Eine Kosten-Nutzen-Analyse unter
    Einbeziehung von Fertigungs-, Material- und Systemintegrationskosten ist f\"ur
    jede der vorgestellten Topologien ein eigenst\"andiger Forschungsgegenstand.
\end{itemize}

%%--------------------------------------------------------------
\abschnitt{Kennzahlentabelle: Quantitative Zusammenfassung}
%%--------------------------------------------------------------

Tabelle~\ref{tab:kennzahlen} fasst die wichtigsten quantitativen Ergebnisse aller
Kapitel in einer komprimierten Form zusammen. Sie dient als schnelle Referenz
f\"ur Leser, die gezielt nach Kenngr\"o\ss{}en suchen.

\begin{longtable}{p{3.0cm}p{4.3cm}p{5.2cm}}
\toprule
\textbf{Kapitel / Thema} & \textbf{Wichtigste Kenngr\"o\ss{}e} & \textbf{Zahlenwert / Ergebnis} \\
\midrule
\endhead
\midrule
\multicolumn{3}{r}{\small Fortsetzung n\"achste Seite}\\
\endfoot
\bottomrule
\endlastfoot
Kap.\ 1: MHz-Magnetik &
  Dimensionale Resonanzfrequenz MnZn-Ferrit ($d=5$\,mm) &
  $f_\mathrm{res} \approx 1{,}9$\,MHz \\
Kap.\ 1: MHz-Magnetik &
  Wirbelstromreduktion durch $N$-fache Segmentierung &
  Faktor $1/N$ \\
Kap.\ 2: H$_2$-LLC &
  Nennleistung pro Modul &
  22\,kW \\
Kap.\ 2: H$_2$-LLC &
  Maximaler Systemwirkungsgrad &
  $\eta \approx 98{,}5$\,\% \\
Kap.\ 3: NeMo.bil &
  Ladeleistung Pro-Bus &
  22\,kW (bidirektional) \\
Kap.\ 3: NeMo.bil &
  OBC-Wirkungsgrad &
  97\,\% @ 11\,kW \\
Kap.\ 4: DC-DC EV &
  Ein\-gangs\-span\-nungs\-be\-reich &
  250--370\,V \\
Kap.\ 4: DC-DC EV &
  Maximaler Wirkungsgrad &
  $\eta \approx 98{,}2$\,\% bei 500\,W \\
Kap.\ 4: DC-DC EV &
  Steinmetz-Exponent $\alpha$ &
  1{,}24 (kalorimetrisch) \\
Kap.\ 5: Bidi.\ DC-DC &
  Schaltfrequenz (GaN) &
  200\,kHz \\
Kap.\ 5: Bidi.\ DC-DC &
  Leistungsdichte &
  $> 2{,}1$\,kW/dm$^3$ \\
Kap.\ 6: SOPWM &
  Schaltverluste STPWM &
  289\,W \\
Kap.\ 6: SOPWM &
  Schaltverluste SOPWM (PvSWmin) &
  195\,W ($-32$\,\%) \\
Kap.\ 7: PMSM &
  Teststandsleistung &
  200\,kW (r\"ucksp.\ DC-Versorgung) \\
Kap.\ 7: PMSM &
  Maximale Drehzahl &
  12.000\,min$^{-1}$ \\
Kap.\ 8: ZCS-LLC &
  Nennleistung &
  2\,kW \\
Kap.\ 8: ZCS-LLC &
  Wirkungsgrad Vollbr\"ucke &
  $\eta > 90$\,\% \\
Kap.\ 8: ZCS-LLC &
  Komponenten $L_r$, $C_r$, $L_m$ &
  44{,}6\,$\mu$H, 56{,}8\,nF, 89{,}1\,$\mu$H \\
Kap.\ 9: GaNius &
  THD Eingangsstrom (mit PFC) &
  $< 2$\,\% \\
Kap.\ 9: GaNius &
  Gleichtaktspannung &
  $u_{CM} = 0$ (analytisch) \\
Kap.\ 10: DPT &
  Grenzfrequenz Koaxial-Shunt &
  $f_c \approx 1{,}6$\,MHz \\
Kap.\ 10: DPT &
  Maximalstrom Testaufbau &
  120\,A bei 800\,V \\
Kap.\ 11: LEA-AI &
  RL Konvergenz-Episode &
  $\approx 200$ Episoden \\
Kap.\ 11: LEA-AI &
  DSP-Inferenzzeit &
  $\approx 61$\,$\mu$s (400\,MFLOPS) \\
Kap.\ 12: Microgrid &
  Emulierbare Gesamtleistung &
  4\,MVA \\
Kap.\ 13: NODE &
  RMS-Modellfehler &
  $< 2$\,\% \\
Kap.\ 14: OpenAutoMoLE &
  Optimierungsverfahren &
  NSGA-II, $\mathcal{O}(k|\mathcal{P}|^2)$ pro Gen. \\
\end{longtable}
\addtocounter{table}{-1}
\begin{table}[H]
  \caption{Quantitative Kenngr\"o\ss{}en aller 14 Kapitel im \"Uberblick.}
  \label{tab:kennzahlen}
\end{table}

%%==============================================================
\chapter{Ausblick}
\label{kap:ausblick}
%%==============================================================

\abschnitt{Vorbemerkung: Warum ein ausf\"uhrlicher Ausblick notwendig ist}

Wissenschaftliche Kompendien neigen dazu, den Ausblick als knappe Aufz\"ahlung knapper
Forschungsdesiderate zu behandeln. Dieses Kapitel verfolgt einen anderen Anspruch.
Die Leistungselektronik befindet sich an einem strukturellen Wendepunkt: Die gleichzeitige
Reife von Breitbandl\"ucken-Halbleitern (WBG), tief lernender k\"unstlicher Intelligenz und
der aufkommenden Wasserstoffwirtschaft erzeugt eine technologische Konvergenz, die weit
\"uber inkrementelle Verbesserungen hinausgeht. Dieser Ausblick versucht, die wichtigsten
Entwicklungslinien bis 2040 systematisch zu verfolgen -- von der Bauelementtechnologie
\"uber Systemarchitekturen und Regelungsverfahren bis hin zu gesellschaftlichen und
normativen Rahmenbedingungen.

%%--------------------------------------------------------------
\abschnitt{Halbleitertechnologie: Der Weg jenseits von GaN und SiC}
%%--------------------------------------------------------------

\unterabschnitt{Galliumoxid (Ga$_2$O$_3$) als n\"achste Halbleitergener\-ation}

W\"ahrend GaN und SiC die aktuelle Industrie dominieren, zeichnet sich am Horizont eine
weitere Generation von Breitbandl\"ucken-Halbleitern ab. Galliumoxid (Ga$_2$O$_3$) besitzt
eine Bandl\"ucke von 4{,}5--4{,}9\,eV -- deutlich gr\"o\ss{}er als GaN (3{,}4\,eV) und SiC
(3{,}26\,eV). Dies erm\"oglicht theoretische Durchbruchfeldst\"arken von $\varepsilon_{br}
\approx 8\,\mathrm{MV/cm}$, was zu einem Baliga's Figure of Merit (BFOM) f\"uhrt, das GaN
um mehr als eine Gr\"o\ss{}enordnung \"ubertrifft:
\[
  \mathrm{BFOM}(\mathrm{Ga_2O_3}) = \varepsilon_r \mu_n \varepsilon_{br}^3
  \approx 3444 \cdot \mathrm{BFOM}(\mathrm{Si}).
\]
Die Herstellbarkeit von Ga$_2$O$_3$-Wafern durch Fl\"ussigphasenepitaxie (Czochralski-
Verfahren) -- analog zu Silizium -- k\"onnte die Produktionskosten erheblich senken.
Erste Schottky-Barriere-Dioden mit Sperrspannungen \"uber 3\,kV wurden bereits demonstriert.
Die Hauptherausforderungen bis zur Serienreife liegen in der geringen thermischen
Leitf\"ahigkeit ($\lambda \approx 10\,\mathrm{W/(m\,K)}$ gegen\"uber SiC:
$\lambda \approx 370\,\mathrm{W/(m\,K)}$) und der noch fehlenden p-Dotierung.

\unterabschnitt{Diamant als ultimativer Leistungshalbleiter}

Kubischer Diamant besitzt mit Bandl\"ucke 5{,}47\,eV, thermischer Leitf\"ahigkeit
$\lambda = 2000\,\mathrm{W/(m\,K)}$ und Durchbruchfeldst\"arke $\varepsilon_{br} =
10\,\mathrm{MV/cm}$ theoretisch die besten Eigenschaften aller bekannten Halbleiter f\"ur
Leistungsanwendungen. Erste p-Kanal-Diamant-MOSFETs wurden an der KEK-Forschungseinrichtung
(Japan) und am Diamond Electronics Lab (Ohio State) demonstriert. Der Weg zur Serienreife
erfordert die L\"osung fundamentaler Probleme: die Herstellung ausreichend gro\ss{}er
Einkristalle mit kontrollierbarer Dotierung (Bor f\"ur p-Typ, Phosphor f\"ur n-Typ),
stabile Gateoxide und reproduzierbare Epitaxieprozesse. Mit einem Zeithorizont von
15--25 Jahren erscheinen erste Nischenprodukte (Hochspannungsdioden, HF-Transistoren)
realistisch.

\unterabschnitt{GaN: Die n\"achste Frequenzgrenze und Monolithische Integration}

Im Bereich GaN wird die Schaltfrequenzgrenze in den n\"achsten zehn Jahren von derzeit
1--10\,MHz in Richtung 30--100\,MHz verschoben. Dies erfordert:

\begin{itemize}
  \item \textbf{GaN-on-Diamond-Substrat}: Der Einsatz von polykristallinem Diamant als
    W\"armesenke unter der GaN-Epitaxieschicht erlaubt die Abfuhr der deutlich erh\"ohten
    Verlustleistungsdichten bei MHz-Schaltfrequenzen. Erste kommerzielle Produkte sind
    f\"ur den milit\"arischen und Raumfahrtbereich zug\"anglich; der kommerzielle Markt
    wird in den 2030er-Jahren erwartet.
  \item \textbf{Monolithische GaN-ICs}: Die Integration von Gate-Treiber, Schutzschaltung,
    Stromsensor und Leistungstransistor auf einem einzigen GaN-Chip (GaN-on-Si-IC)
    eliminiert parasit\"are Schaltschleifen-Induktivit\"aten nahezu vollst\"andig und
    erm\"oglicht Totzeiten unter 1\,ns.
  \item \textbf{E-Mode GaN mit verbesserter Zuverl\"assigkeit}: Aktuelle Normierungsaktivit\"aten
    (JEDEC JEP182) arbeiten an standardisierten Qualifizierungsverfahren f\"ur GaN-HEMTs,
    die den Weg zu AEC-Q101 und ISO-26262-konformen Komponenten ebnen.
  \item \textbf{Bidirektionale GaN-Schalter (wie in GaNius)}: Monolithisch bidirektionale
    GaN-Transistoren, die in \autoref{kap:ganius} vorgestellt wurden, werden von der
    Forschung in die breite Anwendung \"ubergehen. Die JEDEC-Standardisierung solcher
    Bauelemente ist f\"ur 2027--2029 geplant.
\end{itemize}

%%--------------------------------------------------------------
\abschnitt{Magnetik: Vom Ferrit zur integrierten Spule}
%%--------------------------------------------------------------

\unterabschnitt{Nano- und Mikrokompositwerkstoffe f\"ur den MHz-Bereich}

Die in \autoref{kap:mhz_magnetik} untersuchten Grenzen der Ferritechnologie --
insbesondere die dimensionale Resonanz -- werden durch neue Werkstoffklassen adressiert:

\begin{itemize}
  \item \textbf{Nanokristalline Legierungen (Vitroperm, Finemet)}: Durch Nanokristallisierung
    amorpher Metallb\"ander entstehen Werkstoffe mit hoher S\"attigungsfluss\-dichte
    ($B_{sat} \approx 1{,}2\,\mathrm{T}$), hoher Permeabilit\"at ($\mu_r \approx 20.000$)
    bei gleichzeitig niedrigen Verlusten bis 500\,kHz. Gerollte Ringkerne aus
    Vitroperm 800\,F eignen sich als Ersatz f\"ur MnZn-Ferrite bis in den MHz-Bereich.
  \item \textbf{Metallpulververbundwerkstoffe (Powder Core)}: Durch Isolierung einzelner
    Metallk\"orner ($d_k \approx 1\,\mu\mathrm{m}$) in einem Polymerbinder entstehen
    Materialien mit verteiltem Luftspalt und DC-Bias-Festigkeit bei Frequenzen bis
    10\,MHz. Aktuelle Forschung bei Micrometals und Arnold Magnetic Technologies
    zielt auf $\mu_r \approx 50$--200 bei $< 100$\,kHz Verlust-Grenzfrequenz.
  \item \textbf{Ferrit-Nanokomposit-Hybride}: Die gezielte Einlagerung von NiZn-Ferrit-
    Nanopartikeln in eine Polymer-Matrix erlaubt die Abstimmung von $\mu_r$ und
    $\varepsilon_r$ unabh\"angig voneinander -- damit kann die dimensionale Resonanz-
    frequenz nach \autoref{satz:dimres} gezielt auf Werte weit oberhalb der Schaltfrequenz
    verschoben werden.
\end{itemize}

\unterabschnitt{On-Chip-Induktoren und 3D-Integration}

Mit steigender Schaltfrequenz ($f_s > 10\,\mathrm{MHz}$) werden die ben\"otigten
Induktivit\"atswerte so klein, dass eine Integration auf dem Halbleiterchip oder im
Package realisierbar wird:
\begin{itemize}
  \item \textbf{Air-Core-Spulen} in CMOS-Technologie (Kupfer-Metallisierung, 3--7 Lagen)
    erreichen $L = 1$--$50\,\mathrm{nH}$ mit G\"utefaktoren $Q = 10$--$30$ bei 100\,MHz.
  \item \textbf{Embedded Magnetics}: Das Einlaminieren von Ferritmaterial in FR4-Leiterplatten
    (IPower-Technologie, TDK) erlaubt L-Werte von 10--500\,nH direkt im PCB-Substrat.
  \item \textbf{3D-Stacked-Power-Module}: Vollst\"andige Leistungswandler (Transistoren,
    Treiber, Spulen, Kondensatoren) in einem $< 1\,\mathrm{cm}^3$-Geh\"ause mit
    Leistungsdichten $> 10\,\mathrm{kW/dm^3}$ sind f\"ur die 2030er-Jahre prognostiziert.
\end{itemize}

%%--------------------------------------------------------------
\abschnitt{Regelungsverfahren: Von klassischer Regelungstechnik zur erkl\"arbaren KI}
%%--------------------------------------------------------------

\unterabschnitt{Die Grenzen klassischer Kaskadenregelung}

Die in \autoref{kap:pmsm} beschriebene feldorientierte Regelung mit PI-Kaskadenstruktur
erreicht bei hochdynamischen Antrieben (Einstellzeit $< 100\,\mu\mathrm{s}$) und stark
nichtlinearen Betriebsbereichen (Feldschw\"ache, Ges\"attigung) ihre Grenzen. Drei
Entwicklungsrichtungen sind zu erwarten:

\paragraph{Modellpr\"adiktive Regelung (MPC) f\"ur Leistungselektronik.}
Direkte MPC (DMPC) optimiert in jedem Abtastschritt den Spannungsraumzeiger durch
Minimierung einer Kostenfunktion \"uber einen Pr\"adiktionshorizont $N_p$:
\[
  \min_{u_k,\ldots,u_{k+N_p-1}} \sum_{j=0}^{N_p-1}
  \Bigl[\|x_{k+j} - x_{ref}\|_Q^2 + \|u_{k+j}\|_R^2\Bigr],
\]
mit Zustandsgewicht $Q \succ 0$ und Stellgr\"o\ss{}engewicht $R \succ 0$.
Gegen\"uber PI-Reglern erlaubt MPC die explizite Ber\"ucksichtigung von Strom-,
Spannungs- und Drehmomenten-Beschr\"ankungen in der Optimierung.
Der Berechnungsaufwand bleibt durch Explicit-MPC (Vorausberechnung aller aktiven
Mengen offline) oder durch effiziente Solver (OSQP, qpOASES) in Echtzeit handhabbar.

\paragraph{Erkl\"arbare k\"unstliche Intelligenz (XAI) f\"ur Regelungsentscheidungen.}
Ein zentrales Hindernis beim Einsatz von RL-Reglern (vgl.\ \autoref{kap:rl}) in
sicherheitskritischen Anwendungen ist die fehlende Interpretierbarkeit der gelernten Policy.
Methoden der erkl\"arbaren KI -- insbesondere SHAP (SHapley Additive exPlanations) und
LIME (Local Interpretable Model-agnostic Explanations) -- k\"onnen post-hoc erkl\"aren,
welche Zustandsvariablen die Regelentscheidung beeinflusst haben. Die Integration solcher
Erkl\"arungen in den Zertifizierungsprozess nach ISO~26262 und ISO~21448 (SOTIF) ist
ein aktives Forschungsfeld.

\paragraph{Foundation Models f\"ur Leistungselektronik.}
Gro\ss{}e vortrainierte Zeitreihenmodelle (z.B.\ TimeGPT, Chronos) k\"onnen durch
wenige Messsequenzen auf neue Konverterto\-po\-logien adaptiert werden (Few-Shot Learning).
Ein solches Foundation Model k\"onnte die LEA-AI-Forschung (\autoref{kap:rl}) und
ML-Expert (\autoref{kap:ml_expert}) zu einem einheitlichen Rahmenwerk zusammenf\"uhren:
ein Modell, das sowohl Systemidentifikation als auch Regelungsstrategie liefert.
Voraussetzung ist ein gro\ss{}er, \"offentlicher Datensatz aus realen Antriebssystemen --
ein Analogon zu ImageNet f\"ur Leistungselektronik.

\unterabschnitt{Sicherheitszertifizierung lernbasierter Regler}

Die rechtliche und normative Zulassung lernbasierter Regler in Automobilen (ASIL~D)
und Energieanlagen (IEC~61508 SIL~3) ist die gr\"o\ss{}te nichtechnische Herausforderung.
Folgende Ans\"atze werden verfolgt:

\begin{itemize}
  \item \textbf{Control Barrier Functions (CBF)}: Mathematische Garantien, dass der
    Zustand des geregelten Systems eine vorgegebene Sicherheitsmenge $\mathcal{S}$
    nie verl\"asst:
    \[
      \dot{h}(x, u) \geq -\alpha(h(x)) \quad \forall x, u,
    \]
    wobei $h : \mathbb{R}^n \to \mathbb{R}$ die Sicherheitsfunktion ($\mathcal{S} = \{x \mid h(x) \geq 0\}$)
    und $\alpha$ eine Funktion der Klasse $\mathcal{K}$ ist. CBF-basierte Regler k\"onnen
    jeden beliebigen nominellen Regler (also auch einen RL-Agenten) um eine Sicherheitsh\"ulle
    erg\"anzen.
  \item \textbf{Formale Verifikation durch neuronale Lyapunov-Funktionen}: Neuronale Netze
    k\"onnen als Lyapunov-Funktionskandidaten $V(x)$ trainiert werden, sodass
    $V(x) > 0$ und $\dot{V}(x) < 0$ auf der gesamten Zustandsraumdomains gelten -- ein
    formaler Stabili\-t\"ats\-nachweis. Tools wie DREAL und Z3 k\"onnen solche Aussagen
    verifizieren.
  \item \textbf{Runtime-Monitoring}: Online-\"Uberwachung des Reglerausgangs durch einen
    parallelen, zertifizierten klassischen Regler; bei Anomalieerkennung \"ubernimmt der
    klassische Regler (Fail-Safe-Modus). Dieses Konzept ist bereits f\"ur autonom fahrende
    Fahrzeuge etabliert und auf Leistungselektronik \"ubertragbar.
\end{itemize}

%%--------------------------------------------------------------
\abschnitt{Wasserstoffwirtschaft: Vom kW- zum TW-Ma\ss{}stab}
%%--------------------------------------------------------------

\unterabschnitt{Die Notwendigkeit der Hochskalierung}

Aktuelle Szenarien zur Erf\"ullung des Pariser Klimaabkommens prognostizieren eine
Glo\-bale Wasserstoff-Erzeugungskapazit\"at von 600--1000\,Mt/a bis 2050, wovon der
gr\"o\ss{}te Teil durch Elektrolyse aus erneuerbarer Energie stammen soll~\cite{IEA2023}.
Ausgehend von einem Energiebedarf von ca.\ 55\,kWh/kg\,H$_2$ ergibt sich ein Leistungsbedarf
f\"ur die Elektrolyse von:
\[
  P_{Elektrolyse} = \frac{800 \cdot 10^9\,\mathrm{kg/a} \cdot 55\,\mathrm{kWh/kg}}
  {8760\,\mathrm{h/a}} \approx 5\,\mathrm{TW}.
\]
Zum Vergleich: die gesamte heutige globale Stromerzeugungskapazit\"at betr\"agt ca.\ 8\,TW.
Die in \autoref{kap:h2} beschriebenen 22-kW-Module m\"ussen also um sechs Gr\"o\ss{}enordnungen
skaliert werden -- ein Skalierungsproblem ohne historisches Vorbild in der Leistungselektronik.

\unterabschnitt{Modulare Mehrpegelumrichter (MMC) f\"ur Gro\ss{}elektrolyse}

Der Modulare Mehrpegelumrichter (MMC) ist die Schl\"usseltechnologie f\"ur den direkten
Mittelspannungsanschluss gro\ss{}er Elektrolyseur-Farmen. Jedes Submodul (SM) des MMC
enth\"alt einen Kondensator und eine H-Br\"ucke; die Serienschaltung von $N$ Submodulen
pro Arm erzeugt eine nahezu sinusf\"ormige Ausgangsspannung:
\[
  u_{out}(t) \approx \sum_{k=1}^{N} u_{SM,k}(t) \approx V_{DC} \cdot \sin(\omega t + \varphi).
\]
F\"ur eine Eingangsspannung von 10\,kV und Submodul-Kapazit\"at von 800\,V sind
$N = 12$--$16$ Submodule pro Arm erforderlich. Die Gleichstr\"ome f\"ur den Elektrolyseur
k\"onnen dann aus den Submodul-Kondensatoren direkt abgezapft werden (hybride MMC-Struktur).
Aktuelle Forschung konzentriert sich auf:
\begin{itemize}
  \item Energieausgleichsregelung zwischen Submodulen bei dynamischer Lastand\"erung
    des Elektrolyseur-Stacks,
  \item Fehlertoleranz: Kurzschluss eines Submoduls soll den Gesamtbetrieb nicht unterbrechen
    (Bypass-Thyristor-Konzept),
  \item Reduktion der Submodul-Kondensatorgr\"o\ss{}e durch pr\"adiktive Energieregelung.
\end{itemize}

\unterabschnitt{Elektrochemie-informierte Konverterregelung}

PEM-Elektrolyseur-Stacks zeigen bei raschen Leistungs\"anderungen (Rampen $> 10\,\%/s$)
transiente \"Uberspannungen durch Ionentransportlimitierungen und Membranhydratation.
Ein rein elektrischer Regler, der nur Spannung und Strom misst, reagiert zu sp\"at.
K\"unftige Regler werden elektrochemische Zustandsmodelle (Protonentransport, Gasblasendynamik,
Membranwiderstand als Funktion von Feuchte und Temperatur) als Pr\"adiktionsmodell in einer
MPC-Struktur nutzen:
\[
  \dot{x}_{chem}(t) = f_{chem}(x_{chem}(t), I_{stack}(t), T_{stack}(t)),
\]
wobei $x_{chem}$ den elektrochemischen Zustand des Stacks beschreibt. Durch Vorausberechnung
der Stackreaktion k\"onnen \"Uberspannungen vermieden und die Lebensdauer der Membran
erh\"oht werden.

\unterabschnitt{Gr\"uner Wasserstoff als Systemdienstleistungserbringer}

Eine weitgehend unbekannte, aber wichtige Perspektive: Gro\ss{}elektrolyseure k\"onnen
durch gezielte Lastregelung Systemdienstleistungen f\"ur das Stromnetz erbringen:
\begin{itemize}
  \item \textbf{Prim\"arregelreserve (PRL)}: Schnelle Leistungsanpassung ($< 30\,\mathrm{s}$)
    zur Frequenzst\"utzung bei Netzst\"orungen. Der Elektrolyseur reduziert oder erh\"oht
    seinen Strombezug proportional zur Frequenzabweichung.
  \item \textbf{Blindleistungskompensation}: Der netzseitige Umrichter des Elektrolyseurs
    kann unabh\"angig von seiner Wirkleistung Blindleistung in das Netz einspeisen oder
    absorbieren -- analog zu einem STATCOM.
  \item \textbf{Trägheits-Emulation (Synthetic Inertia)}: Durch schnelle Strom\"anderungen
    kann der Umrichter eine virtuelle Massentr\"agheit $H_{virt}$ emulieren, die zur
    Stabilisierung von Netzen mit hohem Anteil erneuerbarer Energien ohne Synchrongeneratoren
    beitr\"agt.
\end{itemize}

%%--------------------------------------------------------------
\abschnitt{Elektromobilit\"at: Konvergenz von Fahrzeug, Netz und Infrastruktur}
%%--------------------------------------------------------------

\unterabschnitt{Vehicle-to-Grid als Systemressource im Energienetz}

Das in \autoref{kap:nemo} vorgestellte NeMo.bil-Konzept ist der Ausgangspunkt einer
weitaus gr\"o\ss{}eren \"Okonomie. Bei 50 Millionen E-Fahrzeugen in Europa (Prognose 2040)
mit durchschnittlich 70\,kWh Batteriekapazit\"at und 50\,\% Verf\"ugbarkeit f\"ur V2G
steht eine disponibler Energiespeicher von:
\[
  E_{V2G} = 50 \cdot 10^6 \cdot 70\,\mathrm{kWh} \cdot 0{,}5 = 1{,}75\,\mathrm{TWh}
\]
zur Verf\"ugung. Zum Vergleich: Die gesamte europ\"aische Pumpspeicherkapazit\"at betr\"agt
ca.\ 200\,GWh. V2G k\"onnte damit zur gr\"o\ss{}ten Speichertechnologie Europas werden --
wenn die technischen, regulatorischen und wirtschaftlichen Rahmenbedingungen stimmen.

Die Realisierung setzt voraus:
\begin{itemize}
  \item \textbf{Hardware}: Bidirektionale OBCs mit h\"ochster Effizienz (vgl.\ \autoref{kap:nemo}),
    genormte Kommunikationsschnittstellen (ISO~15118-20) und langlebige Batteriezellen
    (minimale Degradation durch V2G-Zyklen -- eine Forschungsfrage, die Elektrochemiker und
    Leistungselektroniker gemeinsam l\"osen m\"ussen).
  \item \textbf{Regelung}: Aggregations-Algorithmen, die die Leistungszusagen Hunderttausender
    individueller Fahrzeuge zu einem verl\"asslichen Regelenergieprodukt zusammenfassen.
    Hierbei spielen Multi-Agenten-Systeme und dezentrale Optimierung eine Schl\"usselrolle.
  \item \textbf{Regulierung}: Die Ver\"anderung des Akkreditierungsrahmens f\"ur
    Regelenergieanbieter (aktuell: Mindestangebotsgr\"o\ss{}e 1\,MW) auf kleinere
    Einheiten ist notwendige Voraussetzung f\"ur V2G-Aggregatoren.
\end{itemize}

\unterabschnitt{Dynamisches Laden w\"ahrend der Fahrt}

Das drahtlose Laden w\"ahrend der Fahrt (Dynamic Wireless Power Transfer, DWPT) ist die
konsequente Weiterentwicklung des NeMo.bil-Konvoi-Konzepts. Induktive Ladespulen in der
Fahrbahn \"ubertragen Energie an das fahrende Fahrzeug -- prinzipiell ohne Stopp. Die
physikalischen Grundlagen beruhen auf resonanter Induktionskopplung:
\[
  \eta_{WPT} = \frac{k^2 Q_1 Q_2}{(1 + k^2 Q_1 Q_2)^2},
\]
wobei $k$ der Kopplfaktor und $Q_1$, $Q_2$ die G\"utefaktoren von Sende- und Empfangsspule
sind. Bei GaN-basierten Invertern und Resonanzfrequenzen von 85\,kHz (SAE~J2954) bzw.
zuk\"unftig 1\,MHz (f\"ur kleinere Spulen) sind Systemwirkungsgrade von 85--93\,\% erreichbar.
Die Herausforderungen f\"ur die Masseneinf\"uhrung sind:
\begin{itemize}
  \item Kostenreduktion der Fahrbahninfrastruktur (aktuell: $\sim 1$\,M\,EUR/km),
  \item Standardisierung der Frequenz und Spulengeometrie (SAE, IEC, ISO),
  \item Elektromagnetische Vertr\"aglichkeit: Streufelder m\"ussen f\"ur Passanten und
    Schrittmachertr\"ager unbedenklich sein (ICNIRP-Leitlinien),
  \item Regelung der Leistungs\"ubertragung bei variablem Luftspalt und Querversatz.
\end{itemize}

%%--------------------------------------------------------------
\abschnitt{Microgrid, Verteilnetze und der digitale Energieraum}
%%--------------------------------------------------------------

\unterabschnitt{Solid-State-Transformer: Die universelle Netzkoppeleinheit}

Der Solid-State-Transformer (SST) ist die technologische Antwort auf die zunehmende
Komplexit\"at des Energiesystems. In einem einzigen Ger\"at kombiniert er:
\begin{itemize}
  \item galvanische Trennung (wie ein konventioneller Transformator),
  \item aktive Oberschwingungsfilterung (wie ein Aktivfilter),
  \item AC/DC-Umwandlung (wie ein PFC-Gleichrichter),
  \item DC-DC-Wandlung auf mehreren Spannungsebenen (wie ein DC-DC-Konverter),
  \item Reactive Power Support (wie ein STATCOM).
\end{itemize}
Die Leistungselektronik eines SST f\"ur die 10-kV-Mittelspannungsebene ben\"otigt
Bauelemente mit $V_{DSS} \geq 3{,}3\,\mathrm{kV}$ und schnellen Schaltzeiten -- genau
die Nische, in der SiC-MOSFETs heute und Ga$_2$O$_3$-Transistoren morgen ihren gr\"o\ss{}ten
Wert entfalten.

\unterabschnitt{Digitaler Zwilling des Energienetzes}

Das in \autoref{kap:microgrid} vorgestellte Microgrid-Labor l\"asst sich als physischer
Kern eines digitalen Zwillings verstehen. Ein vollst\"andiger digitaler Netz-Zwilling
spiegelt in Echtzeit den Zustand jedes einzelnen Betriebsmittels (Leitungen, Transformatoren,
Umrichter, Speicher):
\[
  \hat{x}_{Netz}(t) = f_{Zwilling}(\hat{x}_{Netz}(t-\Delta t),\, u_{mess}(t),\, \theta_{Modell}),
\]
wobei $\theta_{Modell}$ durch kontinuierliches Lernen an die Realit\"at angepasst wird.
Anwendungen:
\begin{itemize}
  \item \textbf{Pr\"adiktive Netzplanung}: Simulation der Auswirkungen neuer Windparks oder
    Elektrolyseure auf Spannungsprofile und Kurzschlussleistung, bevor die Anlage gebaut wird.
  \item \textbf{Cyber-Sicherheit}: Erkennung von Anomalien im Netzverhalten durch Vergleich
    von digitalem Zwilling und Realsystem -- ein fr\"uhes Warnsignal f\"ur Cyberangriffe.
  \item \textbf{Black-Start-Planung}: Simulation des Wiederanlaufs nach einem vollst\"andigen
    Netzausfall, um die optimale Reihenfolge der Kraftwerks- und Microgrid-Zuschaltung zu
    bestimmen.
\end{itemize}

%%--------------------------------------------------------------
\abschnitt{Messtechnik und Qualit\"atssicherung: Der untersch\"atzte Fortschrittsmotor}
%%--------------------------------------------------------------

\unterabschnitt{Der Double-Pulse Test in der Zukunft}

Das in \autoref{kap:dpt} analysierte DPT-Verfahren wird in den n\"achsten Jahren
grundlegend weiterentwickelt:

\begin{itemize}
  \item \textbf{Automatisierter Roboter-DPT}: Vollautomatische Messstationen, die Schalt\-parameter
    (Spannung, Strom, Temperatur, Gatewiderstand) systematisch variieren und
    ML-basiert auswerten, erlauben die Charakterisierung ganzer Transistor-Wafer in Stunden
    statt Wochen.
  \item \textbf{In-situ-DPT w\"ahrend des Betriebs}: Methoden zur Online-Charakterisierung
    von Transistoren im laufenden Konverterbetrieb (ohne Unterbrechung) erm\"oglichen die
    Erkennung von Alterungsph\"anomenen (NBTI, Gate-Oxiderm\"udung, Bond-Wire-Ersch\"opfung)
    in Echtzeit -- essenziell f\"ur Predictive Maintenance.
  \item \textbf{Breitband-Feld-Sondensysteme}: Das aktuelle Hauptproblem der DPT-Messtechnik
    -- die Einschr\"ankung durch Sondenkapazit\"at und Shunt-Induktivit\"at (vgl.\ \autoref{kap:dpt})
    -- wird durch Entwicklung miniaturisierter Rogowski-Spulen mit $> 2\,\mathrm{GHz}$
    Bandbreite und elektrooptischer Spannungssensoren adressiert.
\end{itemize}

\unterabschnitt{Kalorimetrie: Von der Einzelkomponente zum Gesamtsystem}

Die kalorimetrische Verlustmessung (\autoref{kap:dcdc_ev}) wird k\"unftig auf h\"ohere
Systemebenen ausgedehnt:
\begin{itemize}
  \item \textbf{Systemkalorimetrie}: Messung des Gesamtwirkungsgrads eines vollst\"andigen
    Wandlersystems (inkl.\ Gate-Treiber, Steuerplatine, K\"uhlsystem) durch Einbau in eine
    thermisch isolierte Kammer -- der \glqq{}Goldstandard\grqq{} f\"ur Wirkungsgradangaben.
  \item \textbf{Thermische Bildgebung mit hoher Aufl\"osung}: Moderne InSb-IR-Kameras
    ($< 10\,\mathrm{mK}$ Temperaturaufl\"osung, $> 1\,\mathrm{MHz}$ Framerate) erm\"oglichen
    die r\"aumlich aufgel\"oste Verlustanalyse einzelner Bonddr\"ahte und Chip-Regionen.
\end{itemize}

%%--------------------------------------------------------------
\abschnitt{Systemintegration: Leistungselektronik als Plattformtechnologie}
%%--------------------------------------------------------------

\unterabschnitt{Konvergenz aller behandelten Themen}

Die 14 Kapitel dieses Kompendiums beschreiben keine isolierten Entwicklungen -- sie
sind Komponenten eines integrierten Systems, das in den 2030er-Jahren entstehen wird.
\autoref{abb:systemkonvergenz} skizziert die Wechselwirkungen:

\begin{itemize}
  \item \textbf{GaN/SiC-Halbleiter} (Kap.\ 1) $\rightarrow$ erm\"oglichen MHz-Schaltfrequenzen
    $\rightarrow$ erfordern neue \textbf{Magnetwerkstoffe} (Kap.\ 1) $\rightarrow$ erlauben
    kompaktere \textbf{LLC-Konverter} (Kap.\ 2, 8).
  \item \textbf{Bidirektionale Wandler} (Kap.\ 3, 5) $\rightarrow$ erm\"oglichen \textbf{V2G}
    $\rightarrow$ stabilisieren \textbf{Microgrids} (Kap.\ 13) $\rightarrow$ integrieren
    \textbf{Elektrolyseure} (Kap.\ 2).
  \item \textbf{RL-Regler} (Kap.\ 11) + \textbf{NODE-Modelle} (Kap.\ 13) $\rightarrow$
    bilden zusammen \textbf{OpenAutoMoLE} (Kap.\ 14) der n\"achsten Generation: ein
    selbstoptimierendes, datengetriebenes Auslegungswerkzeug.
  \item \textbf{DPT-Messtechnik} (Kap.\ 10) $\rightarrow$ liefert Schaltenergie-Daten
    $\rightarrow$ trainiert \textbf{Schaltverlustmodelle} f\"ur SOPWM (Kap.\ 6) und
    OpenAutoMoLE (Kap.\ 14).
\end{itemize}

\unterabschnitt{Standardisierung und Normgebung als Enabler}

Technologische Exzellenz allein gen\"ugt nicht -- ohne Normierung und Standardisierung
bleiben Innovationen im Labor. Die wichtigsten anstehenden Normierungsaktivit\"aten:

\begin{itemize}
  \item \textbf{IEC~62684 / ISO~15118-20 (V2G-Kommunikation)}: Die endg\"ultige Ratifizierung
    von ISO~15118-20 (Bidirektionales Laden, AC und DC) ist der rechtliche Schl\"ussel
    zur Skalierung von V2G-Anwendungen. Technisch sind alle Voraussetzungen erf\"ullt;
    die Normierung schafft die notwendige Interoperabilit\"at.
  \item \textbf{IEC~60747-15 (SiC-Bauelementpr\"ufung)}: Standardisierte Pr\"ufverfahren
    f\"ur SiC-MOSFETs, die \"uber den Stand des DPT hinausgehen (Langzeitpr\"ufungen,
    Cosmic-Ray-Tests bei Hochspannung, Body-Diode-Charakterisierung).
  \item \textbf{IEC~62909 (Bidirektionale Netzwechselrichter)}: Normierung der netz-
    schutzrelevanten Eigenschaften bidirektionaler Wechselrichter (Inselerkennung,
    Synchronisierung, FRT -- Fault Ride-Through).
  \item \textbf{ISO~21448 (SOTIF) f\"ur KI-Regler in Leistungselektronik}: Die Safety
    of the Intended Functionality-Norm, urspr\"unglich f\"ur autonomes Fahren entwickelt,
    wird auf lernbasierte Regler in der Antriebstechnik ausgedehnt.
\end{itemize}

%%--------------------------------------------------------------
\abschnitt{Offene Forschungsfragen und Zukunftsaufgaben}
%%--------------------------------------------------------------

Abschlie\ss{}end werden die wichtigsten offenen Forschungsfragen systematisch benannt.
Sie sind Einladungen an die n\"achste Generation von Leistungselektronik-Forschern.

\begin{aufgabe}[Optimalit\"at simultaner Mehrfachziele bei SOPWM]
  Kann ein Pulsmuster gleichzeitig THD minimieren, Schaltverluste minimieren und
  ZVS-Bedingungen gew\"ahrleisten? Oder existiert ein fundamentaler Trade-off, der
  durch einen Satz von Pareto-optimalen Pulsmuster-Familien beschrieben werden muss?
  Ist das kombinierte Optimierungsproblem konvex oder NP-hart?
\end{aufgabe}

\begin{aufgabe}[Globale Stabilit\"at lernbasierter Regler]
  Unter welchen Bedingungen kann ein durch RL gelernter Regler f\"ur einen Leistungskonverter
  mit einer globalen Lyapunov-Funktion zertifiziert werden? Gibt es eine untere Schranke
  f\"ur die Gr\"o\ss{}e des neuronalen Netzes, unterhalb derer keine globale Stabilit\"at
  m\"oglich ist?
\end{aufgabe}

\begin{aufgabe}[Thermodynamische Wirkungsgradgrenze f\"ur LLC-Konverter]
  Was ist der theoretisch maximal erreichbare Wirkungsgrad f\"ur einen 22-kW-LLC-Konverter
  bei gegebenen Bauelementtechnologien (GaN, Vitroperm-Kern, Kupferwicklung)?
  Ist die aktuelle Technik ($\eta \approx 98{,}5\,\%$) bereits nahe an diesem Limit,
  oder sind durch bessere Topologien und Materialien noch mehrere Zehntel Prozent zu gewinnen?
\end{aufgabe}

\begin{aufgabe}[Skalierungsgesetze f\"ur Microgrid-Stabilit\"at]
  Gibt es Phasen\"uberg\"ange in der Stabilit\"at eines Microgrids, die auftreten, wenn die
  Zahl der Teilnehmer $N$ einen kritischen Wert \"uberschreitet? Lassen sich Methoden der
  statistischen Physik (Mean-Field-Theorie, Renormierungsgruppe) auf die Netzwerktopologie
  anwenden?
\end{aufgabe}

\begin{aufgabe}[Batterie-Degradation durch V2G-Zyklen]
  Wie stark beeinflusst der Einsatz einer Fahrzeugbatterie als V2G-Speicher ihre Lebensdauer?
  Welche Ladestrategie (Spannungsgrenzen, Stromstufen, Ruhepausen) maximiert die Nutzbarkeit
  der Batterie f\"ur V2G bei gleichzeitig minimalem Kapazit\"atsverlust?
\end{aufgabe}

\begin{aufgabe}[Koaxial-Shunt-Alterung und Langzeitcharakterisierung]
  Wie ver\"andert sich die Impedanzkennlinie eines Koaxial-Shunts \"uber seine Lebensdauer
  (1000\,h, 5000\,h Betrieb bei Hochstrom) und bei Temperaturwechselbelastung? Gibt es
  Vorhersagemodelle f\"ur den Zeitpunkt der Kalibrierungspflicht?
\end{aufgabe}

\begin{aufgabe}[Interaktion von KI-Regler und Netzschutz]
  Wenn ein RL-Regler die Regelstrategie dynamisch anpasst, kann er unvorhergesehene
  Wechselwirkungen mit dem Netzschutz (\"Uberstromausl\"oser, Distanzschutz) produzieren.
  Wie sind solche Interaktionen systematisch zu modellieren und zu verhindern?
\end{aufgabe}

\begin{aufgabe}[Quantencomputing f\"ur Leistungselektronik-Optimierung]
  In welchem Zeithorizont kann Quantencomputing praktische Vorteile f\"ur die
  multikriterielle Optimierung (OpenAutoMoLE) bieten? Bei welcher Problemgr\"o\ss{}e
  (Anzahl Entwurfsvariablen, Zielfunktionen) ist ein Quanten-NSGA-II einem klassischen
  NSGA-II \"uberlegen?
\end{aufgabe}

\medskip

Die Leistungselektronik steht am Beginn ihrer wichtigsten Dekade. Die Herausforderungen --
Klimaschutz, Energiewende, Mobilit\"atstransformation -- sind gro\ss{}, aber die Werkzeuge,
die dieses Kompendium beschreibt, sind real und einsatzbereit. Welche Antworten die
Forschungsgemeinschaft auf die offenen Fragen dieses Ausblicks findet, wird ma\ss{}geblich
bestimmen, ob die Energiewende gelingt.


%%==============================================================
%% Schriftenverzeichnis
%%==============================================================
\begin{thebibliography}{99}

\bibitem{IEA2023}
International Energy Agency (IEA),
\emph{Global Hydrogen Review 2023},
IEA, Paris, 2023. \url{https://www.iea.org/reports/global-hydrogen-review-2023}

\bibitem{Asgarnia2023}
R.~Asgarnia,
\emph{GaNius -- Neue gleichstaktarme dreiphasige PFC-Gleichrichterfamilie
mit bidirektionalen GaN-Transistoren},
Forschungsbericht, LEA Universit\"at Paderborn, 2023.

\bibitem{Baliga2021}
B.\,J.~Baliga,
\emph{Fundamentals of Power Semiconductor Devices}, 2.~Auflage,
Springer, Cham, 2021.

\bibitem{Blaabjerg2017}
F.~Blaabjerg, Y.~Yang, D.~Yang, X.~Wang,
``Distributed Power-Generation Systems and Protection,''
\emph{Proc.\ IEEE}, Bd.~105, Nr.~7, S.~1311--1331, Juli 2017.

\bibitem{Bolte2020}
S.~Bolte, N.~Fr\"ohleke, J.~B\"ocker,
``DC-DC Converter Design for Power Distribution Systems in Electric Vehicles
using Calorimetric Loss Measurements,''
\emph{Proc.\ EPE-ECCE Europe}, 2020.

\bibitem{Buttler2018}
A.~Buttler, H.~Spliethoff,
``Current Status of Water Electrolysis for Energy Storage, Grid Balancing
and Sector Coupling via Power-to-Gas and Power-to-Liquids,''
\emph{Renewable and Sustainable Energy Reviews}, Bd.~82,
S.~2440--2454, 2018.

\bibitem{Chen2018}
T.\,Q.~Chen, Y.~Rubanova, J.~Bettencourt, D.\,K.~Duvenaud,
``Neural Ordinary Differential Equations,''
\emph{Advances in Neural Information Processing Systems (NeurIPS)}, 2018.

\bibitem{Deb2002}
K.~Deb, A.~Pratap, S.~Agarwal, T.~Meyarivan,
``A Fast and Elitist Multiobjective Genetic Algorithm: NSGA-II,''
\emph{IEEE Trans.\ Evol.\ Comput.}, Bd.~6, Nr.~2, S.~182--197, Apr.~2002.

\bibitem{Dixon2001}
L.\,H.~Dixon,
\emph{Ferrite Core Loss for Power Conversion},
Unitrode Power Supply Design Seminar, Texas Instruments, 2001.

\bibitem{Foerster2021}
N.~F\"orster, D.~Urbaneck, B.~Kohlhepp, D.~K\"ubrich, O.~Wallscheid, F.~Schafmeister,
``Pitfalls and their Avoidability in the Double-Pulse Test,''
\emph{Proc.\ IEEE ECCE}, 2021.

\bibitem{Foerster2023}
N.~F\"orster,
\emph{OpenAutoMoLE -- Open-Source-Software zur automatisierten Modellierung
leistungselektronischer Konverter},
Proc.\ ETG Fachtagung, 2023.

\bibitem{Forrester2008}
A.\,I.\,J.~Forrester, A.~S\'obester, A.\,J.~Keane,
\emph{Engineering Design via Surrogate Modelling},
Wiley, Chichester, 2008.

\bibitem{Grieves2017}
M.~Grieves, J.~Vickers,
``Digital Twin: Mitigating Unpredictable, Undesirable Emergent Behavior
in Complex Systems,''
in \emph{Transdisciplinary Perspectives on Complex Systems},
Springer, S.~85--113, 2017.

\bibitem{jones2020gan}
E.\,A.~Jones, F.\,F.~Wang, D.~Costinett,
``Review of Commercial GaN Power Devices and GaN-Based Converter
Design Challenges,''
\emph{IEEE Journal of Emerging and Selected Topics in Power Electronics},
Bd.~4, Nr.~3, S.~707--719, Sept.~2016.

\bibitem{Keuck2021}
L.~Keuck,
\emph{Bidirektionaler DC-DC-Wandler f\"ur Mehrspannungs-Bordnetze
in Automobilanwendungen},
Studentenprojekt, LEA Universit\"at Paderborn, 2021.

\bibitem{Kurs2007}
A.~Kurs, A.~Karalis, R.~Moffatt, J.\,D.~Joannopoulos,
P.~Fisher, M.~Soljacic,
``Wireless Power Transfer via Strongly Coupled Magnetic Resonances,''
\emph{Science}, Bd.~317, Nr.~5834, S.~83--86, Juli 2007.

\bibitem{Mengelkamp2018}
E.~Mengelkamp, J.~G\"arttner, K.~Rock, S.~Kessler, L.~Orsini, C.~Weinhardt,
``Designing Microgrid Energy Markets: A Case Study,''
\emph{Applied Energy}, Bd.~210, S.~870--880, 2018.

\bibitem{Morimoto1994}
S.~Morimoto, Y.~Takeda, T.~Hirasa, K.~Taniguchi,
``Expansion of Operating Limits for Permanent Magnet Motor by Current
Vector Control Considering Inverter Capacity,''
\emph{IEEE Trans.\ Ind.\ Appl.}, Bd.~26, Nr.~5, S.~866--871, 1990.

\bibitem{Muhlethaler2012}
J.~M\"uhlethaler, J.~Biela, J.\,W.~Kolar, A.~Ecklebe,
``Core Losses under the DC Bias Condition Based on Steinmetz Parameters,''
\emph{IEEE Trans.\ Power Electron.}, Bd.~27, Nr.~2, S.~953--963, Feb.~2012.

\bibitem{NemoBil2024}
J.~Wiegard et al.,
\emph{NeMo.bil -- Netzbasiertes Mobiles Laden im \"offentlichen Nahverkehr:
Systemkonzept und Realisierung},
Forschungsbericht, LEA Universit\"at Paderborn, 2024.

\bibitem{Peter2022}
K.~Peter, S.~Hanke, J.~B\"ocker, F.~Mink,
``Inverter Loss Management for an Electrical High-Speed Drive System,''
\emph{Proc.\ IEEE ECCE Asia}, 2022.

\bibitem{Raissi2019}
M.~Raissi, P.~Perdikaris, G.\,E.~Karniadakis,
``Physics-Informed Neural Networks: A Deep Learning Framework for Solving
Forward and Inverse Problems Involving Nonlinear PDEs,''
\emph{J.\ Computational Physics}, Bd.~378, S.~686--707, 2019.

\bibitem{Schulman2017}
J.~Schulman, F.~Wolski, P.~Dhariwal, A.~Radford, O.~Klimov,
``Proximal Policy Optimization Algorithms,''
\emph{arXiv preprint arXiv:1707.06347}, 2017.

\bibitem{Sutton2018}
R.\,S.~Sutton, A.\,G.~Barto,
\emph{Reinforcement Learning: An Introduction},
2.~Auflage, MIT Press, Cambridge, 2018.

\bibitem{Urbaneck2022}
D.~Urbaneck, J.~Wiegard, F.~Schafmeister,
``Analysis of Inverter Operation Modes of an IGBT-Based ZCS LLC Converter
for a 2\,kW Automotive On-Board DC-DC,''
\emph{Proc.\ IEEE APEC}, 2022.

\bibitem{Weber2022}
D.~Weber, D.~Schmies,
\emph{Microgrid-Labor -- Emulation dezentraler Energienetze mit 4\,MVA},
Technischer Bericht, LEA Universit\"at Paderborn, 2022.

\end{thebibliography}


%%==============================================================
\chapter{Abbildungsverzeichnis}
%%==============================================================
\listoffigures


\end{document}
